\section{Completion, finite coefficients, and the arithmetic action}
\label{sec:E}

We work in characteristic zero.  In this section $\overline{\mathcal M}_{g,n}$
denotes the stable compactification, while
$\mathcal M_{g,n,\overline{\mathbb Q}}$ denotes geometric base change of
the smooth moduli stack.  This distinction is fixed by the maps in
\S\ref{sec:E-manuscript}: the bar in cote 143 is geometric base change.
Write $\Gamma_{g,n}$ for the orientation-preserving mapping class group
fixing the $n$ labels, and $C_{g,n}$ for a finite-object skeleton of its
boundary exit category.  Arrows point from a nodal stratum towards a
smoothing.  All actions and homomorphisms between profinite groups below
are continuous.

\subsection{Finite congruences and the comparison with sheaves}

\begin{result}{E.1}{Proved}{The strict completion and its finite representations}
For every stable $(g,n)$, set
\[
 \widehat C_{g,n}=\varprojlim_R C_{g,n}/R,
\]
where $R$ ranges over finite-index congruences on all Hom-sets, with the
object set fixed.  Composition makes this a category internal to
profinite sets, and restriction along $C_{g,n}\longrightarrow\widehat
C_{g,n}$ is an equivalence
\[
 \operatorname{Fun}_{\mathrm{cts}}(\widehat C_{g,n},\mathrm{FinSet})
       \simeq \operatorname{Fun}(C_{g,n},\mathrm{FinSet}).       \tag{E.1}
\]
Every functor between two such finite-object categories extends to their
strict completions.  A natural transformation extends as well, and the
extensions preserve its equations.
\end{result}
\begin{proof}
There are finitely many stable labelled dual graphs of a fixed type:
the number of edges is at most $3g-3+n$, the number of vertices is at
most $2g-2+n$, and every vertex genus and every incidence set is bounded.
Jansen's Construction 3.18 identifies isomorphism classes of exit objects
with these graphs.  Choose one object in each class.

For a finite congruence $R$, the quotient keeps identities and all
composition laws.  Intersection of two finite congruences is finite,
because its classes inject into the product of their finite class sets.
The inverse system is therefore cofiltered.  Taking inverse limits of
the finite composition functions defines continuous composition, and
associativity follows in every finite coordinate.

For $F:C_{g,n}\to\mathrm{FinSet}$, equality of the functions $F(a)$ and
$F(b)$ defines a finite congruence on each Hom-set.  There are finitely
many objects, so only finitely many finite function sets occur.  Thus
$F$ factors through one quotient.  Conversely, a continuous map from a
profinite Hom-set to a finite function set factors through finitely many
coordinates; a common refinement handles all pairs of objects.  All
transition maps and projections are surjective on Hom-sets, so the
identity and composition equations descend to this common finite stage.
The components of a natural transformation are finitely many functions;
its naturality equations can be checked on the dense original Hom-sets
or at a common finite stage.  This proves full faithfulness and essential
surjectivity in (E.1).

For a functor $u:C\to D$, pull back every finite congruence on $D$ along
$u$.  The resulting finite congruences on $C$ define the extension
$\widehat u$.  This construction respects composition of functors.
Its application to the finitely many components of a natural
transformation gives the last assertion.  No identification of
$\widehat\Gamma_{g,n}$ with its procongruence image occurs in this proof.
\end{proof}



\begin{result}{E.2}{Known}{Constructible finite coefficients in every stable type}
Fix $\overline{\mathbb Q}\hookrightarrow\mathbb C$.  For every stable
$(g,n)$ there are equivalences, compatible with stalk functors,
\[
 \operatorname{Cons}_{\partial,\mathrm{fin}}
       ((\overline{\mathcal M}_{g,n})_{\overline{\mathbb Q},\mathrm{et}})
 \simeq
 \operatorname{Cons}_{\partial,\mathrm{fin}}
       ((\overline{\mathcal M}_{g,n})^{\mathrm{an}})
 \simeq
 \operatorname{Fun}_{\mathrm{cts}}(\widehat C_{g,n},\mathrm{FinSet}).
                                                               \tag{E.2}
\]
Here $\operatorname{Cons}_{\partial,\mathrm{fin}}$ means finite-stalk
sheaves locally constant on every topological-type stratum of the fixed
boundary stratification.  Jansen, Corollary 6.6 and Theorem 6.7, and
Barwick--Glasman--Haine, Proposition 12.6.4, give the inputs.
\end{result}
\begin{proof}[Derivation, including the stack descent]
The compactification is a finite-type separated Deligne--Mumford stack.
Choose a finite-type quasi-compact etale atlas $U\to\overline{\mathcal
M}_{g,n}$.  Its iterated fibre products form a simplicial finite-type
scheme after refinement by affine etale covers.  On each such scheme,
Proposition 12.6.4 of \cite{ET_BGH2020} identifies constructible etale
sheaves with constructible analytic sheaves.  The comparison is inverse
image, so its compatibility with every face, degeneracy, and refinement
map is part of the construction.  Descent data and their cocycle
equations are consequently identified.  Taking the descent limit gives
the asserted equivalence for the stack; this is also the construction
of \cite[Proposition 13.9.17; Corollary 13.9.19]{ET_BGH2020}.
Restriction to a stratum commutes with inverse image.  Local constancy
and finiteness are etale-local properties.  Hence the equivalence
restricts to exactly the fixed boundary-constructible subcategories.

For extension from $\overline{\mathbb Q}$ to $\mathbb C$, use invariance
of finite etale covers of finite-type schemes under extension of
algebraically closed characteristic-zero fields.  Apply it on every
stratum and on the normal-crossings punctured local charts; the latter
have finite monodromy described by finite quotients of their meridian
lattices.  The specialization maps are maps of finite sets commuting
with these monodromies.  The same invariance thus identifies the
gluing data, and etale descent gives the stack assertion.

For the analytic term, Jansen's Corollary 6.6 identifies the exit
infinity-category with the ordinary Charney--Lee category, in the
variance of Construction 3.18.  Her Theorem 6.7 supplies exodromy.
Take set-valued objects, then the full subcategory with finite values.
Passing to the chosen skeleton is an equivalence; (E.1) completes the
derivation.  We have used the stable-type case of Corollary 6.6 directly,
so no unstable-type inequality printed in a subsequent specialization
is needed.
\end{proof}

\begin{result}{E.3}{Refuted}{The proposed nonsolvable-centralizer obstruction to (E.2)}
The claim in Part E(a) that the finite constructible-sheaf comparison
requires
$C_{\widehat\Gamma}(H)=\overline{C_\Gamma(H)}$ for the Klein-quartic
isotropy group $H\cong\operatorname{PSL}_2(\mathbb F_7)$ is false as
a claimed necessary step.  Formula (E.2) holds there, with this isotropy
retained, by its descent proof.  It does not assert the displayed
centralizer identity.
\end{result}
\begin{proof}
In an etale local quotient chart $[V/H]$, a finite constructible sheaf
is an $H$-equivariant finite constructible sheaf on $V$: the equivariance
maps are indexed by the actual finite set $H$ and satisfy its
multiplication table.  Analytification preserves those maps and their
equations, and its inverse on constructible sheaves transports them
back.  This proves the equivalence on $[V/H]$ for every finite $H$,
including the indicated group of order $168$.  The descent limit
used in E.2 then glues these comparisons.  Neither functor in this
explicit comparison takes a centralizer in a completed mapping class
group.  The surviving action of $H$ on the stalk over a fixed point
is precisely its stack isotropy action.  This proves the positive
comparison at the point singled out in the prompt and removes the
asserted dependency by an explicit construction.
\end{proof}

\begin{result}{E.4}{Proved}{Functorial arithmetic action and its faithfulness}
Forgetting, clutching, and relabelling, with their specified
compatibility isomorphisms, induce functors on the strict completions.
The full completed system has a coherent continuous outer action of
$G_{\mathbb Q}=\operatorname{Gal}(\overline{\mathbb Q}/\mathbb Q)$,
and the map to equivalence classes of its autoequivalences is injective.
At a rational normal-crossings tangential point, a compatible local
representative acts on the meridian lattice by
\[
 (t_e)_{e\in E(G)}\longmapsto
                 (t_e^{\chi(\sigma)})_{e\in E(G)} .           \tag{E.3}
\]
Changing geometric paths conjugates the simultaneous local action.
\end{result}
\begin{proof}
Apply E.1 to the discrete geometric functors and their coherence maps
constructed in Parts A and C.  Alternatively, use (E.2): the algebraic
maps are defined over $\mathbb Q$, and inverse image preserves finite
boundary-constructible sheaves because each source stratum maps into
one target stratum.

For $\sigma\in G_{\mathbb Q}$, the functor $(\sigma^{-1})^*$ on
geometric sheaves gives a left action.  Composition of inverse image
supplies its coherent associativity.  Every finite cover, finite stalk,
and finite gluing diagram used by a constructible sheaf descends to a
number field, so this action is continuous on finite data.  Transferring
through (E.2), including its stalk functors, gives the outer categorical
action.  Compatibility with the geometric operations follows from
their being defined over $\mathbb Q$.

For a smoothing parameter $q_e$, the covers $q_e^{1/m}$ have inertia
$\mu_m$.  The action on these roots of unity is $\chi\bmod m$.
Taking all $m$ proves (E.3), with a single exponent on the whole
local lattice.  A changed path conjugates the inertia subgroup.

For faithfulness, restrict the system to the open stratum of
$\overline{\mathcal M}_{0,4}$, namely
$\mathbb P^1-\{0,1,\infty\}$, with group $\widehat F_2$.
An outer-trivial arithmetic action fixes every conjugacy class of open
subgroups, hence every connected finite cover of the tripod and its
smooth compactification.  Belyi's theorem gives such a presentation
for every elliptic curve over $\overline{\mathbb Q}$.  The action
therefore fixes every elliptic $j$-invariant.  Every algebraic number
is an elliptic $j$-invariant, so the Galois element is the identity.
This argument concerns the full system.  Its $(0,3)$ labelled component
alone is terminal and has the trivial arithmetic action.
\end{proof}

\subsection{The five-point higher comparison}

\begin{result}{E.5}{Proved}{Primitive meridians and retracts at all five-point flags}
Let $G=\pi_1(\mathcal M_{0,5})$.  For a multicurve $\sigma$ on the
labelled five-punctured sphere put $H_\sigma=\operatorname{Stab}_G(\sigma)$.
Every $H_\sigma$ is a retract of $G$.  For $|\sigma|=1$ it is
$\mathbb Z\times F_2$; for $|\sigma|=2$ it is $\mathbb Z^2$.
Every full flag of boundary strata has a connected lifting groupoid,
whose automorphism group is one of
\[
 G,\quad F_2,\quad \mathbb Z\times F_2,\quad
 \mathbb Z^2,\quad\mathbb Z,\quad 1.                         \tag{E.4}
\]
All these groups are good, and the finite categorical quotients induce
their full profinite topologies.
\end{result}
\begin{proof}
Normalize labels $1,2,3$ to $0,1,\infty$, and write the remaining
coordinates as $x,y$.  The open moduli space is the complement of
the five lines
$x=0,x=1,y=0,y=1,x=y$ in $\mathbb C^2$.
Their positive meridians form a basis of $H_1(G,\mathbb Z)=\mathbb Z^5$:
small transverse discs give the generators, and integration of the
five logarithmic differentials gives the dual integral homomorphisms.
The ten boundary meridians have the following columns in this basis:
\[
\begin{array}{c|rrrrr}
12&-1&-1&-1&-1&-1\\
13&0&1&0&1&1\\
14&1&0&0&0&0\\
15&0&0&1&0&0\\
23&1&0&1&0&1\\
24&0&1&0&0&0\\
25&0&0&0&1&0\\
34&-1&-1&0&0&-1\\
35&0&0&-1&-1&-1\\
45&0&0&0&0&1
\end{array}                                                   \tag{E.5}
\]
These are valuations of $x,x-1,y,y-1,x-y$ along the boundary.
For example, at $D_{12}$ put $x=q^{-1}$ and $y=a q^{-1}$ with
$a\notin\{0,1\}$; all five valuations are $-1$.
At $D_{13}$ take $x=1+q,y=1+aq$, and at $D_{23}$ take
$x=q,y=aq$.  The two remaining infinity divisors give the rows
$34,35$ directly.

For a single curve $c$ enclosing a labelled pair $\{a,b\}$,
the cutting sequence is
\[
1\longrightarrow\langle t_c\rangle\longrightarrow H_c
 \xrightarrow{r}\Gamma_{0,4}\longrightarrow1.               \tag{E.6}
\]
Here $r$ is the restriction of the actual forgetting homomorphism
$G\to\Gamma_{0,4}$ obtained by forgetting $a$; on the normalized
boundary it replaces the node by the remaining mark $b$.
Each row in (E.5) is primitive, so there is an integral character
$\ell:G\to\mathbb Z$ with $\ell(t_c)=1$.
The restriction of $(\ell,r):G\to\mathbb Z\times F_2$ to $H_c$
is injective: its kernel lies in $\langle t_c\rangle$ by (E.6),
and $\ell$ kills no nonidentity power.  It is surjective: lift the
$F_2$ coordinate using (E.6) and correct the integer coordinate by
a power of $t_c$.  Its inverse on $H_c$, composed with $(\ell,r)$,
is the required retraction $G\to H_c$.

A two-curve stratum has all three normalized components of type
$(0,3)$, so its stabilizer is the free lattice generated by its two
twists.  Its labels are disjoint pairs, giving exactly fifteen
possibilities.  Every corresponding pair of rows of (E.5) has a
$2\times2$ minor of determinant $1$ or $-1$.  The program
\texttt{code/E\_meridian\_pairs.py} checks all fifteen and writes an
explicit integral left inverse for each two-column matrix.  Composing
this left inverse with $G\to H_1(G,\mathbb Z)$ gives a retraction
onto that twist lattice.  Conjugating the chosen multicurve conjugates
the retraction, so this proves the assertion for every representative.

Consider now a strict flag represented, after individual choices of
transporters, by
$\sigma_0\supsetneq\cdots\supsetneq\sigma_r$.
Pure labelled sphere mapping classes fix each edge of a stabilized
graph, so $H_{\sigma_0}\subset H_{\sigma_i}$ for every $i$.
Successive changes of target trivialization bring every lifted flag to
the canonical one.  Naturality of an automorphism along its arrows
then gives exactly
\[
             \operatorname{Aut}(\text{flag})
                =H_{\sigma_0}/\Delta_{\sigma_r}.             \tag{E.7}
\]
There is no assumption that all representatives for all flags can be
chosen simultaneously as inclusions.  Formula (E.7) is invariant under
the transporters used for one flag.  The possible cardinalities are
$0,1,2$, and (E.6) gives exactly (E.4).

A retract inherits the full profinite topology: any finite quotient of
the retract extends along its retraction.  The indicated lattice
quotients also inherit their quotient profinite topology.  For goodness,
the projection $(x,y)\mapsto x$ is the configuration fibration with
base $\mathbb C-\{0,1\}$ and fibre
$\mathbb C-\{0,1,x\}$, so
$G=F_3\rtimes F_2$ (the group extension splits because its quotient
is free).  Finitely generated free groups and free abelian groups are
good.  For completeness, their free or Koszul resolutions compare
ordinary and continuous finite-coefficient cohomology.  In the free
extension, characteristic finite-index subgroups of $F_3$ are cofinal;
the action on each finite quotient has finite image.  This proves the
exactness of the completed split extension and continuity of the
finite cohomology modules.  Comparing its two Lyndon--Hochschild--Serre
spectral sequences proves goodness of $G$.  The same argument proves
goodness for the direct products in (E.4).
\end{proof}

\begin{result}{E.6}{Proved}{LM3b H.9 at five points}
The nerve of the strict finite-congruence completion of $C_{0,5}$
represents the profinite stratified homotopy type of
$\overline{\mathcal M}_{0,5}$ with its boundary stratification.
In particular it classifies the corresponding constructible sheaves
with $\pi$-finite values, with all higher gluing retained.
\end{result}
\begin{proof}
We verify all full flags, including point--divisor--open flags.
Every component of a genus-zero multicurve determines a split of
the puncture labels.  Distinct disjoint components have distinct
splits: two disjoint curves giving the same split cobound a
puncture-free annulus and are isotopic.  A pure mapping class
preserves each split.  Thus, for fixed representatives $\sigma,\tau$,
every transporter $g$ with $g\sigma\supseteq\tau$ has the same
subset $g^{-1}\tau\subseteq\sigma$.  Consequently
\[
 \{g\in G:g\sigma\supseteq\tau\}=H_\tau g_0
\]
for any one transporter $g_0$.  This is transitivity on the actual
incidence set, not merely on abstract graph types.  It also proves
$H_{\sigma_0}\subseteq H_{\sigma_i}$ for an actual nested flag:
a pure stabilizer fixes each component individually.  The same
argument on the cut components shows that each flag of fixed
labelled graph types has one $G$-orbit.  The labelled splits specify
the retained edges at every step, and homeomorphisms of the
complementary punctured spheres identify any two realizations.

For $N\lhd G$ of finite index define the finite quotient category
using the images in $G/N$ of the transporter sets and of the source
twist subgroups.  These categories are cofinal among the finite
congruence quotients.  Indeed, a functor to a finite category has a
finite-index kernel on each target automorphism group.  Pull it back
to $H_\tau$.  By E.5 there is $N_\tau\lhd G$ of finite index whose
intersection with $H_\tau$ lies in this kernel.  Intersect over the
finitely many targets.  Two transporters with equal image modulo
$N\Delta_\sigma$ differ, after a source twist, by an element of
$N\cap H_\tau$: the preceding fixed-incidence calculation shows
that their ratio fixes $\tau$.  They therefore have equal functor
images.  This
constructs the desired factorization.

For a flag as in (E.7), the lifting groupoid at this finite stage is
transitive, with automorphism group
\[
 p_N(H_{\sigma_0})/p_N(\Delta_{\sigma_r}).                    \tag{E.8}
\]
This follows by the same successive target changes of trivialization.
Every finite-stage arrow lifts to an original transporter arrow
between the chosen objects, so a finite-stage chain has simultaneous
composable lifts.  The transitivity just proved passes to this
lifting groupoid.  We do not assume that an arbitrary quotient of
a curve complex retains all its incidence information.
The nested inclusions $H_{\sigma_0}\subset H_{\sigma_i}$ ensure that
the final automorphism lies in $p_N(H_{\sigma_0})$; its ambiguity is
exactly $p_N(\Delta_{\sigma_r})$.  Thus no extra intersection or
fibre-product assertion is being inferred from a single crossing.

By E.5, system (E.8) is cofinal in all finite quotients of (E.7),
and its group is good.  Goodness for every finite local system,
together with the identical profinite fundamental group, proves that
the corresponding system of classifying spaces is the profinite
homotopy completion of the original flag groupoid.  One may verify
this last implication by obstruction theory into a finite Postnikov
tower: each successive obstruction and choice is a finite local-system
cohomology group, for which goodness is precisely the required
comparison.

All these maps are induced by restriction of actual chains, so they
commute with every face map.  The spatial decollage theorem
\cite[Theorem 2.7.4; Construction 2.8.8 and (2.8.9)]{ET_BGH2020}
reconstructs the profinite stratified object from this full flag
diagram.  The comparison is an equivalence at every entry, hence is
an equivalence of the reconstructed objects.  Finally Corollary
12.6.7 of that source identifies the latter object with the
stratified etale type.  We have checked the three-stratum flags
explicitly; the rank-one argument of LM3b H.8 alone would not give
this step.
\end{proof}

\subsection{The arithmetic subgroup selected by a real framing}

\begin{result}{E.7}{Known}{The AA selection in ANABELIAN E.14}
Fix the standard embedding $G_{\mathbb Q}\hookrightarrow\widehat{
\operatorname{GT}}$ and its complex conjugation $c=(-1,1)$.
Among the maximal closed subgroups satisfying the AA-property of
Hoshi--Mochizuki--Tsujimura, Definition 5.12, exactly one contains $c$;
it is the specified arithmetic image.  This statement retains all
prerequisites of that definition, including its configuration-space
compatibility.  It confirms the subgroup-selection assertion of
ANABELIAN E.14.
\end{result}
\begin{proof}[Derivation from the opened sources]
Theorem 5.17(ii) of \cite{ET_HMT2025} identifies the maximal AA
subgroups with the conjugates of the arithmetic image $G$; its
proof uses Theorem 5.15(iv), Remark 5.12.1, and Lemma 5.16.
Proposition 4(ii) of \cite{ET_LS1997} gives
$C_{\widehat{\operatorname{GT}}}(c)=\langle c\rangle$.
We checked its proof on printed p.~581: commutation forces the
second GT coordinate to be fixed by simultaneous inversion, its
fixed subgroup is trivial, and the hexagon gives multiplier $\pm1$.

Let $bGb^{-1}$ contain $c$.  Then $b^{-1}cb$ is a nontrivial
involution in $G$.  All such involutions in $G_{\mathbb Q}$ are
conjugate: their fixed fields are real closures, the field $\mathbb Q$
has a unique ordering, and an isomorphism of real closures extends
to their algebraic closures.  Thus $b^{-1}cb=g c g^{-1}$ for some
$g\in G$, whence $bg\in C_{\widehat{\operatorname{GT}}}(c)
\subset G$.  It follows that $b\in G$ and $bGb^{-1}=G$.
Existence follows because the arithmetic image itself satisfies AA.

This is a test on closed subgroups of the completed genus-zero
configuration system; Hoshi--Mochizuki--Tsujimura construct it from
the abstract group of the second tripod configuration space.  It
does not identify every individual element satisfying inertia and
permutation conditions with a Galois element.  Once the subgroup
has been selected, membership in it is the additional test.
\end{proof}

\subsection{A mixed-braid calculation in genus one}

\begin{result}{E.8}{Proved}{The kernel for an automorphism of an extension}
Let $1\to K\to E\to Q\to1$ be an exact sequence of finitely generated
profinite groups, with $Z(K)=Z(Q)=1$, and let
$\rho:Q\to\operatorname{Out}(K)$ be its outer monodromy.  Restriction
and passage to the quotient give an isomorphism
\[
 \ker\bigl(\operatorname{Out}(E;K)\longrightarrow
                    \operatorname{Out}(Q)\bigr)
       \simeq C_{\operatorname{Out}(K)}(\rho(Q)).             \tag{E.8}
\]
The same formula holds after imposing any invariant peripheral
condition on the restriction to $K$, by intersecting the right side
with the corresponding subgroup of $\operatorname{Out}(K)$.
\end{result}
\begin{proof}
There is a canonical isomorphism of profinite groups
\[
 E\longrightarrow
 \{(q,a)\in Q\times\operatorname{Aut}(K):[a]=\rho(q)\},
 \qquad e\longmapsto (\bar e,\operatorname{conj}_e|_K).
                                                               \tag{E.9}
\]
The kernel is $Z(K)$.  To prove surjectivity, lift $q$ to $e\in E$;
then $a(\operatorname{conj}_e|_K)^{-1}$ is inner, so multiplying $e$
by an appropriate element of $K$ gives the prescribed pair.  The
continuous bijection is an isomorphism, since the groups are compact
and Hausdorff.  This proof does not require injectivity of $\rho$.

An automorphism of $E$ inducing the identity on $Q$ restricts to an
outer automorphism centralizing $\rho(Q)$.  Conversely, for a lift
$a\in\operatorname{Aut}(K)$ of such an outer automorphism, the map
$(q,b)\mapsto(q,aba^{-1})$ of the right side of (E.9) defines an
automorphism of $E$ over $Q$.  Changing $a$ by an inner automorphism
changes this lift by an element of $\operatorname{Inn}(K)$.
If an automorphism over $Q$ is inner in $E$, its conjugating element
maps to $Z(Q)=1$, and hence belongs to $K$.  Finally, an element of
the kernel on the left of (E.8) can be modified by an inner
automorphism of $E$ to induce the identity on $Q$.  These operations
are inverse and prove the assertion, including its peripheral version.
\end{proof}

\begin{result}{E.9}{Proved}{An explicit inertia-preserving GT lift to
the three-point elliptic group}
Put $K=\widehat\pi_{1,2}$, $Q=\widehat\Gamma_{1,2}$ and
$E=\widehat\Gamma_{1,3}$, where all marks are ordered and the
map $E\to Q$ forgets the third mark.  The standard GT action on
$\widehat B_5$ induces a continuous homomorphism
\[
 \Psi:\widehat{\operatorname{GT}}\longrightarrow
           \operatorname{Out}(E;K)
                                                               \tag{E.10}
\]
whose quotient action on $Q=\widehat B_4/Z(\widehat B_4)$ is the
standard action of Minamide--Nakamura, Theorem~4.3(i).  Its restriction
to $K$ fixes both cusp classes with the common multiplier $\lambda$,
and its action on $E$ preserves nonseparating twist inertia.
\end{result}
\begin{proof}
We specify the quotient, its topology, and the action.  Let
$B_{4,1}\subset B_5$ be the subgroup whose permutation fixes strand
$5$.  Forgetting strand $5$ splits by the first four strands and gives
\[
 B_{4,1}=F_4\rtimes B_4,
 \quad F_4=\langle x_1,x_2,x_3,x_4\rangle,
 \quad x_i=x_{i5}.
\]
Set $c=x_1x_2x_3x_4$ and $\omega_4=(\sigma_1\sigma_2\sigma_3)^4$.
Our Artin convention is
\[
 \sigma_i(x_i)=x_{i+1},\qquad
 \sigma_i(x_{i+1})=x_{i+1}^{-1}x_ix_{i+1},
\]
with the other generators fixed.  Thus every $\sigma_i$ fixes $c$,
and $\omega_4$ acts by $\operatorname{Inn}(c^{-1})$.  These identities
follow by successively applying the displayed formulas; the exact
reduced-word checks for four, and also three, generators are recorded
in \texttt{code/E\_braid\_models.py} and its output.
The actual identification $x_i=x_{i5}$, including
this conjugation direction, is checked in the
faithful Artin action of $B_3,B_4,B_5$ by
\texttt{E\_normalizer\_checks.py}; that program also
checks the first-strand cabling used in E.28.

Complete this split sequence.  Its kernel has the full profinite
topology: characteristic finite-index subgroups of $F_4$ are cofinal,
and the action on each corresponding finite quotient factors through
a finite quotient of $B_4$.  Hence
$\widehat B_{4,1}=\widehat F_4\rtimes\widehat B_4$.
Write
\[
 D=\widehat F_4/\overline{\langle\!\langle x_i^2:1\leq i\leq4
                              \rangle\!\rangle}
       =C_2\amalg C_2\amalg C_2\amalg C_2,
 \qquad K_4=\ker(D\to C_2),                               \tag{E.11}
\]
where every factor maps nontrivially.  We use the same letters for
the images of $x_i,c$.  The discrete index-two covering is the double
cover of a disc branched at four points.  Its surface has genus one
and two boundary components.  Replacing those boundaries by punctures
identifies $K_4$ with $K$; the two peripheral classes are generated by
$c$ and $x_1cx_1^{-1}$.  The Schreier basis $x_1x_i$, $2\leq i\leq4$,
also proves directly that $K_4$ is free profinite of rank three.

Lifting a half-twist of two adjacent branch points to this double
cover gives a Dehn twist on a nonseparating curve.  Thus the Artin
action on $K_4$ is the based version of elliptic universal monodromy.
After removing base points its quotient is the familiar geometric
isomorphism
\[
 B_4/\langle\omega_4\rangle\simeq\Gamma_{1,2};
\]
this is precisely the isomorphism used in
\cite[Corollary C and the preceding paragraph]{ET_MN2022}.
The full profinite center of $\widehat B_4$ is
$\overline{\langle\omega_4\rangle}\simeq\widehat{\mathbb Z}$
\cite[\S2, equations (2.2) and (2.9)]{ET_MN2022}.
Consequently the element $(c,\omega_4)$ is central in
$K_4\rtimes\widehat B_4$ and the quotient
\[
 E'=(K_4\rtimes\widehat B_4)\big/
       \overline{\langle(c,\omega_4)\rangle}          \tag{E.12}
\]
fits into $1\to K_4\to E'\to Q\to1$.  Its central subgroup has
trivial intersection with $K_4$, as is seen by projection to
$\widehat B_4$.  Its outer monodromy is the universal monodromy just
identified.  The geometric universal curve gives the exact sequence
$1\to K\to E\to Q\to1$; this is also the exact sequence in
\cite[Theorem 6.13]{ET_HMTopicsI2012} with all primes retained.
Applying the explicit isomorphism (E.9) to both extensions gives an
isomorphism $E'\simeq E$ inducing the specified identity on kernel and
quotient.  In particular, we have identified the extension itself,
and have not inferred its class from an abstract isomorphism of
the two endpoint groups.

Let $F=(\lambda,f)\in\widehat{\operatorname{GT}}$.  Its genuine
braid action is
\[
 \sigma_1\mapsto\sigma_1^\lambda,
 \qquad
 \sigma_i\mapsto
 f(y_i,\sigma_i^2)^{-1}\sigma_i^\lambda f(y_i,\sigma_i^2),
 \quad y_i=\sigma_{i-1}\cdots\sigma_1^2\cdots\sigma_{i-1}.
                                                               \tag{E.13}
\]
These are \cite[\S0.1, equation (1)]{ET_HS2000}.  They preserve
the permutation quotient and the subgroup on the first four
strands.  They also preserve $\widehat F_4$: the generators $x_{i5}$
are taken to conjugates of their $\lambda$-powers, with conjugators
in the pure subgroup of $\widehat B_5$.  That subgroup lies in
$\widehat B_{4,1}$, where the normal closure of these generators is
$\widehat F_4$.  Applying the inverse action gives equality.
To check conjugacy inside $\widehat F_4$, write such a conjugator
as $vb$, with $v\in\widehat F_4$ and $b\in\widehat P_4$.
The pure Artin action of $b$ sends $x_i$ to an $\widehat F_4$-conjugate
of $x_i$.  This proves the required internal conjugacy as well.

It follows that the closed normal subgroup generated by $x_i^2$
is preserved.  Since $\lambda$ is odd, the parity homomorphism in
(E.11) is preserved too.  The appendix of \cite{ET_HS2000},
printed pp.~34--35, proves the genuine identities
$F(y_i)=y_i^\lambda$ and $F(\omega_i)=\omega_i^\lambda$.
Here $c=y_5$, so both $c$ and $\omega_4$ are taken to their
$\lambda$-powers.  Their commuting quotient relation in (E.12)
is consequently preserved, with
$(c,\omega_4)\mapsto(c,\omega_4)^\lambda$.
We obtain an action on $E'$, hence (E.10), preserving every equation
of the extension.  On its quotient $Q$ the action is exactly (E.13)
for four strands, which proves the asserted comparison map.

The first cusp generator $c$ is taken exactly to $c^\lambda$;
the second is conjugated by an odd element of $D$, and thus remains
in the second $K_4$-conjugacy class with exponent $\lambda$.
Finally, the element $\sigma_1$ in (E.12) is the lift of the
half-twist supported on a small two-branch-point disc.  Its annular
double cover avoids the extra base point, so this lift is a
nonseparating Dehn twist on the three-point torus.  Equation (E.13)
takes it to its $\lambda$-power.  All nonseparating curves are
conjugate under the pure mapping class group of this surface:
cutting along either gives a connected genus-zero surface with the
same three ordered punctures and two ordered boundary components,
and an orientation-preserving identification glues back.  Thus all
their completed inertia subgroups are preserved.  This finishes
the construction without using the high-genus extension argument
of Boggi's Lemma~6.16.
\end{proof}

\begin{result}{E.10}{Proved}{Automatic genus-one inertia: LM3b E.9}
Let $A_{\rm per}$ be the subgroup of
$\operatorname{Out}(\widehat\Gamma_{1,3})$ preserving the kernel
$K=\widehat\pi_{1,2}$ of forgetting the third mark and permuting its
two peripheral subgroup conjugacy classes.  Then
\[
 A_{\rm per}=\langle\delta\rangle\times
                    \Psi(\widehat{\operatorname{GT}})
       \simeq C_2\times\widehat{\operatorname{GT}},          \tag{E.14}
\]
where $\delta$ interchanges the two original labels.  Every element
has one common peripheral exponent, normalizes $\langle\delta\rangle$,
and preserves nonseparating twist inertia.  Thus the group $A$ and
its subgroup $A^I$ in LM3b E.1/E.9 coincide, and
$A/\mathfrak S_2\simeq\widehat{\operatorname{GT}}$.
\end{result}
\begin{proof}
The quotient group $Q=\widehat B_4/Z(\widehat B_4)$ is center-free
\cite[\S2, equation (2.9)]{ET_MN2022}; $K$ is free profinite of
rank three and is center-free.  By (E.8), the kernel of
$A_{\rm per}\to\operatorname{Out}(Q)$ is
\[
 C_{\operatorname{Out}^{C}(K)}(\rho(Q)).                    \tag{E.15}
\]
Here $\operatorname{Out}^{C}$ means permutation of the two
peripheral subgroup classes, without a prescribed exponent.
\cite[Theorem 6.13(ii), equivalently Theorem D(i)]{ET_HMTopicsI2012}
identifies (E.15) with the automorphism group of the tautological
twice-punctured elliptic curve.  Its hypotheses are satisfied:
characteristic zero, all primes, $(g,r)=(1,2)$, $2g-2+r=2>1$,
and the open subgroup there is the whole $Q$.  The theorem and its
proof were checked: for $r>0$ its claims $(\ast1)$ and $(\ast2)$
reduce to Theorem~6.12(i); that theorem reduces to the two-point
configuration case and Lemmas~6.8 and~6.10, and no centralizer
identity for a finite nonsolvable subgroup is used.
The printed induction in Lemma~6.8, p.~132, says $3g-3+r>1$
after the complexity-zero base case.  The argument works for
$3g-3+r>0$.  At the omitted type $(1,1)$ choose the irreducible
one-node stable curve, whose normalization is a tripod.  Inserting
the moving point at its node gives two tripod vertices joined by
two edges.  The two one-edge restrictions in that proof are the
already established tripod case; they kill the two independent
Dehn coefficients, respectively.  At $(0,4)$ use two tripods
joined at a node and the same moving-node insertion.  This supplies
both complexity-one cases, so the application here does not
depend on the erroneous strict inequality.

This tautological automorphism group is $C_2$.  Its nonidentity
element is $z\mapsto p_1+p_2-z$ on the elliptic curve: it exchanges
the two punctures.  In (E.11) its outer action is conjugation by
$x_1$, acting on the Schreier basis $x_1x_i$ by inversion.  Hence
it is nontrivial.  Conjugation by $x_1$ also normalizes
$K_4\rtimes\widehat B_4$ inside $D\rtimes\widehat B_4$ and
preserves the central relation (E.12), so it supplies the same
kernel element in $A_{\rm per}$.  Geometrically it is the
relabeling $\delta$.  Explicitly, conjugation by $x_1$ fixes that
relation because its $D$-component is
\[
 x_1c\omega_4(x_1^{-1})
   =x_1c(c^{-1}x_1^{-1}c)=c.
\]
Thus deck descent has been checked on the defining central
generator itself.  As an orientation-preserving
relabeling it
preserves nonseparating twist inertia and has peripheral exponent
one.

By \cite[Theorem 4.3(i), Corollary C]{ET_MN2022},
$\operatorname{Out}(Q)=\widehat{\operatorname{GT}}$ via the
four-strand action.  E.9 gives a section of the map from
$A_{\rm per}$ onto this group.  Its kernel has order two and is
therefore central: conjugation acts through
$\operatorname{Aut}(C_2)=1$.  This proves (E.14).
The section and its kernel both preserve the stated inertia,
and their peripheral exponents are respectively $\lambda$ and
$1$.  The final assertions follow for every element of the direct
product, including the formerly imposed normalizer condition.
\end{proof}

\subsection{Explicit elliptic boundary comparisons}

\begin{result}{E.11}{Proved}{The two elliptic boundary maps and forgetting}
Put $G_2=\widehat B_4/\langle\omega _4\rangle$ and
$G_1=\widehat B_3/\langle\omega _3^2\rangle$, where
$\omega _r=(\sigma _1\cdots\sigma _{r-1})^r$ and brackets on
central profinite subgroups mean closed subgroups.  Under the
geometric identifications $G_i=\widehat\Gamma_{1,i}$, forgetting
the second mark is
\[
 q:G_2\longrightarrow G_1,\qquad
 (\sigma _1,\sigma _2,\sigma _3)\longmapsto(a,b,a),
 \quad a=\sigma _1,\ b=\sigma _2\text{ in }G_1.       \tag{E.16}
\]
The rational-tail boundary has stabilizer $\widehat B_3\subset G_2$,
twist $\omega _3^2$, and normalization group $G_1$.
For the nonseparating twist $a=\sigma _1\in G_2$, put
$v=\sigma _2\sigma _1^2\sigma _2$ and $d=\sigma _3$.
Its stabilizer $H_a$ is generated by $a,v,d$ and satisfies
\[
 [a,v]=[a,d]=1,\qquad (vd)^2=(dv)^2,\qquad
                    a^2(vd)^2=1.                         \tag{E.17}
\]
Consequently
\[
 H_a/\langle a\rangle\simeq
 \widehat A_3/\langle\omega _3\rangle,
 \quad A_3=\langle\sigma _1^2,\sigma _2\rangle\subset B_3,
                                                               \tag{E.18}
\]
with $\sigma _1^2\mapsto v$ and $\sigma _2\mapsto d$.
This is the fundamental group of
$\mathcal M_{0,\{1,2\},[+,-]}$, the four-point normalization
with its two node branches unordered.  Its geometric branch
character is $v\mapsto1$, $d\mapsto0$ in $C_2$.
For every
$F=(\lambda,f)\in\widehat{\operatorname{GT}}$, the genuine
braid actions commute with (E.16), with the rational-tail
quotient, and with (E.18), including its branch permutation.
\end{result}
\begin{proof}
The geometric braid identification and the Cardano--Ferrari
map are checked in \cite[Corollary C; proof of Proposition
3.2(i), especially (3.1)--(3.5)]{ET_MN2022}.  Keeping the central
involution instead of quotienting it gives precisely (E.16).
The one-boundary torus mapping class group is $B_3$; the
two-chain relation identifies its boundary twist with
$(ab)^6=\omega _3^2$.  Gluing its boundary to the tripod
carrying marks $1,2$ gives the rational-tail stabilizer stated
above.  Its map to $G_2$ is injective: cutting the rational-tail
curve has kernel its infinite twist subgroup and quotient
$\Gamma_{1,1}$, and this is the same central extension as
$B_3\to\mathrm{SL}_2(\mathbb Z)$.

For completeness the centralizer presentation comes directly
from cutting the disk braid along a small disk containing its
first two strands.  A braid preserving that disk is determined
by its internal half twist $a$ and a braid on the collapsed disk
and the other two strands.  The latter braid fixes the collapsed
strand and hence belongs to $A_3$.  The framing of that strand
gives the central $a$-factor.  The two generators of $A_3$ lift
to $v,d$; its presentation is
$\langle v,d\mid vdvd=dvdv\rangle$.
Restoring the four-strand full twist gives
$\omega _4=a^2(vd)^2$ and gives (E.17) after quotienting.
These identities are also verified in the faithful Artin action
on the free group of rank four by the attached program
\texttt{E\_braid\_models.py}.
The cut genus-one surface is a sphere with its original two
punctures and two branches.  The disjoint twist $d$ fixes both
branches.  The order-two element $vd$ modulo $a$ exchanges
them, as the half-turn of the cut cylinder.  Thus the branch
character has values $v\mapsto1,d\mapsto0$.  Its kernel is
the pure four-point group.  In the original mixed braid group
$A_3$ the ordinary strand-permutation character instead has
values $v\mapsto0,d\mapsto1$; these two explicit characters
must be retained in using the quotient identification.

These descriptions pass to the full profinite groups by the
explicit induced-topology calculation in E.26 below.  In
particular, $H_a$ here is the closure of the displayed
geometric stabilizer, with its full profinite topology.
The proof uses the free kernel of the Cardano--Ferrari map
and does not use a profinite twist-centralizer theorem.

The genuine action of $F$ on $B_n$ is
\cite[formula (1)]{ET_HS2000}.  It restricts to the standard
action on $B_3\subset B_4$ and sends $\omega _3$ to
$\omega _3^\lambda$.  This proves the rational-tail assertion.
For (E.16), its only additional generator is $\sigma _3$.
Here $y_3=\sigma _2\sigma _1^2\sigma _2=v$ and
\[
 q(v)=ba^2b=-a^{-2}\quad\hbox{in }\mathrm{SL}_2(\mathbb Z).
\]
The equality follows either from the braid presentation and
$(ab)^6=1$, or from the usual faithful integral matrices
$a=\left(\begin{smallmatrix}1&1\\0&1\end{smallmatrix}\right)$,
$b=\left(\begin{smallmatrix}1&0\\-1&1\end{smallmatrix}\right)$.
Thus $q(v)$ commutes with $a^2$.  Since $f$ lies in the
closed commutator subgroup of the free profinite group,
$f(q(v),a^2)=1$.  Formula (1) therefore sends $q(\sigma _3)$
to $a^\lambda$, exactly as required.  It already commutes on
$\sigma _1,\sigma _2$, proving the equality of genuine maps.

Finally \cite[\S0.3, pp.~7--8]{ET_HS2000} states and derives
the genuine braid-tower action which commutes with the
string-doubling homomorphism
\[
 D:\widehat A_3\longrightarrow\widehat B_4,
 \qquad D(\sigma _1^2)=v,\quad D(\sigma _2)=d.
\]
It also preserves the strand permutation.  Since the actions
send $a$ and $\omega _3$ to their $\lambda$th powers, they
descend through (E.17)--(E.18).  Moreover $v$ maps to $v^\lambda$
and $d$ to a conjugate of $d^\lambda$; $\lambda$ is odd.
They therefore preserve the geometric branch character just
computed, as well as the ordinary strand character.
This proves all the asserted comparisons as equalities before
passing to outer automorphisms.
\end{proof}

\begin{result}{E.12}{Proved}{The orbifold normalizer maps and their parameter}
For $r=3,4$, put $F_r=\widehat F(x_1,\ldots,x_r)$ with its
$r+1$ peripheral classes, $D_r=F_r/\overline{\langle\!\langle
x_1^2,\ldots,x_r^2\rangle\!\rangle}$, and
$K_r=\ker(D_r\to C_2)$, every $x_i$ mapping to the
nonidentity element.  Then $K_3=\widehat\pi_{1,1}$ and
$K_4=\widehat\pi_{1,2}$.  Let $\mathcal A_r$ be the genuine
peripheral automorphism group fixing the infinity class, and
$\widetilde N_r$ the inverse image in $\mathcal A_r$ of the
normalizer of the full geometric sphere group
$\Sigma_{0,r+1}\subset\operatorname{Out}_{\rm lac}(F_r)$.
Quotient-and-restriction
\[
 \rho_r:\mathcal A_r\longrightarrow\operatorname{Aut}_{\rm lac}(K_r)
\]
sends $\widetilde N_r$ into the outer normalizer of the full
elliptic geometric group $\Sigma_{1,r-2}$.  It descends to the
normalizer quotients by these geometric groups and is the
identity on their standard $\widehat{\operatorname{GT}}$
parameters.  The analogous elliptic filling map has the same
parameter.  This settles the orbifold-normalizer portion of
LM3b E.4, with its specified maps.
\end{result}
\begin{proof}
A peripheral automorphism takes $x_i^2$ to a conjugate of
$x_{\sigma(i)}^{2\lambda}$ for a unit $\lambda$.
Applying the same assertion to its inverse proves equality
of the closed normal subgroups generated by the squares.
It preserves parity because $\lambda$ is odd.  With
$c=x_1\cdots x_r$ in $D_r$, the cusp of $K_3$ is $c^2$,
and the two cusps of $K_4$ are $c,x_1cx_1^{-1}$.
An odd conjugator is made even by multiplication by $c$ when
$r=3$, and exchanges the two cusp classes when $r=4$.
This proves the asserted genuine peripheral homomorphism.
Changing its genuine representative by $\operatorname{Inn}(g)$
changes its outer restriction by either $1$ or the geometric
deck involution $\delta_r=[\operatorname{conj}_{x_1}]$, according
to the parity of $g$.

Here is the full normalizer check.  The genus-zero universal
curve gives $\widehat\Gamma_{0,r+2}$ as the inverse image of
the pure geometric subgroup of $\operatorname{Out}(F_r)$ in
$\operatorname{Aut}(F_r)$.  Conjugation therefore embeds its
normalizer quotient into
$\operatorname{Out}(\widehat\Gamma_{0,r+2})$.
Indeed a representative centralizing the inner copy of $F_r$
is the identity, since $F_r$ is center-free.
By \cite[Corollaries 2.6(i), 2.8(i)]{ET_HMM2022}, its target is
$\widehat{\operatorname{GT}}\times\mathfrak S_{r+2}$ with no
extra inertia assumption: here the curve is the tripod and its
configuration length is $r-1\geq2$.  The image normalizes the
geometric point-stabilizer $\mathfrak S_{r+1}$.  Its normalizer
inside $\mathfrak S_{r+2}$ is itself, since its unique fixed
letter is preserved.  Dividing by that geometric subgroup gives
an injection into $\widehat{\operatorname{GT}}$.
Conversely the genuine mixed braid action provides every
parameter.  Thus each element of $\widetilde N_r$ is,
modulo inner automorphisms, a geometric infinity stabilizer
times the restriction $\alpha_F$ of the genuine
$(r+1)$-strand action for a unique $F$.
Both factors preserve infinity, so the geometric factor can
indeed be taken in its stabilizer.

The infinity stabilizer is the image of $B_r$.  Lifting the
branched double cover of its disk identifies its action on
$K_r$ with the elliptic geometric group: for $r=3$ it is
$B_3/\langle\omega _3^2\rangle$, and for $r=4$ it is
$B_4/\langle\omega _4\rangle$ together with $\delta_4$.
The first assertion is the torus two-chain relation; the second
is the braid realization in E.9.  The genuine mixed braid
action preserves $F_r\rtimes\widehat B_r$, intertwining the
action of every $b\in\widehat B_r$ on $K_r$ with that of its
standard $F$-image.  It also normalizes the deck group.
These are all the geometric generators, which proves the full
normalizer inclusion.  It proves descent through changes of
inner representative, infinity stabilizer, and sheet as well.

For $r=4$ the induced action on
$B_4/\langle\omega _4\rangle$ has parameter $F$, precisely
the parameter detected in E.9--E.10.  For $r=3$ the same
mixed-braid construction gives
\[
 E_3=(K_3\rtimes\widehat B_3)/
       \langle(c^2,\omega _3^2)\rangle
       \simeq\widehat\Gamma_{1,2}.                       \tag{E.19}
\]
To verify the extension identification, the branched disk cover
gives its geometric monodromy on $K_3$; the square of the
full twist acts by $\operatorname{Inn}(c^{-2})$.  For a
center-free kernel every extension is the pullback of
$\operatorname{Aut}(K_3)\to\operatorname{Out}(K_3)$ along
this monodromy, by the proof of E.8.  This part of that proof
does not require the quotient to be center-free.
It proves (E.19) and the displayed central relation exactly.
The action sends both entries of that relation to their
$\lambda$th powers and hence descends.

The included $B_3$ in (E.19) consists of the two standard
nonseparating twists on the elliptic component separated by
the rational tail.  It is therefore the stabilizer in E.11,
and its action is the genuine standard $F$-action.  This
detects the parameter in
$\operatorname{Out}(\widehat\Gamma_{1,2})=
\widehat{\operatorname{GT}}$.
Indeed an inner conjugacy in $G_2$ between actions preserving
this $B_3$ normalizes $B_3$ itself.  E.26 proves directly that
its normalizer in $G_2$ is $B_3$.
Restriction thus gives an inner conjugacy in $B_3$ itself.
The standard map
$\widehat{\operatorname{GT}}\to\operatorname{Out}(\widehat B_3)$
is injective: restriction after quotienting its center is the
pair description
$\widehat{\operatorname{GT}}\subset
\widehat{\operatorname{GT}}_0$ of
\cite[\S0.3 and Main Theorem (a)]{ET_HS2000}.
Hence the detected parameter is exactly $F$.
Finally the genuine identity (E.16) proves the same-parameter
assertion for elliptic filling.
\end{proof}

\subsection{The covering-and-filling system}

\begin{result}{E.13}{Proved}{The injectivity asserted in ANABELIAN E.13}
Let $\mathscr H$ be the labelled and coloured category of
ANABELIAN E.8: its objects are oriented hyperbolic surfaces;
its generators are finite unramified covers and specified
puncture fillings; its relations are diagram isotopy and
equality of actual composite maps.  Complete a right mapping-class
orbit $S_u\backslash M_S$ to
$\overline{S_u}\backslash\widehat M_S$ and retain the finite
label-permutation quotients, all label-forgetting maps, and all
specified deletions.  Let $F$ be a continuous type-preserving
autoequivalence preserving this structure.  If its outer action
on every labelled and coloured genus-zero mapping class group
is trivial, its induced outer action on every
$\widehat M_{g,n}$, $2g-2+n>0$, is trivial.

Thus restriction is injective on the compatible families of
outer group automorphisms induced by these autoequivalences.
The statement is about these families, not a classification of
the autoequivalences inducing a given family.  No congruence
subgroup property is used.
\end{result}
\begin{proof}
We give the categorical reconstruction and the required
Hurwitz surjection, since both are essential to the assertion.
There are finitely many right orbits in each uncompleted Hom
set.  A nonidentity elementary arrow strictly decreases
$2g-2+n$, covering degrees are bounded by the ratio of endpoint
complexities, and a finitely generated group has only finitely
many subgroups of bounded index.  Fillings with specified
colours have finitely many orbits.  Composing two finite lists
of orbits still gives finitely many, since the intermediate
mapping class group permutes the first list.  Its stabilizer of
an orbit is consequently open in its full profinite completion.
The completed orbit formula in the statement therefore has
continuous left and right actions and composition; associativity
follows from density.  In particular stabilizers of completed
arrows are exactly closures of the indicated discrete right
stabilizers.  This verifies the feature of E.8 used below.

Put $P_S=\widehat\pi_1(S)$.  Covers give injections of their
specified finite index and fillings give surjections on these
groups, as outer continuous homomorphisms.  These assignments
extend to completed arrows: for finitely generated profinite
groups $P,Q$, continuous homomorphisms $P\to Q$ form a closed
subspace of a finite power of $Q$, and quotienting by target
conjugacy is compact Hausdorff.  Composition is continuous,
as is checked in each finite characteristic quotient.
We use this functor to recognize equations between arrows;
we do not suppose that $F$ already commutes with it.

\emph{Point pushing and chosen sheets.}
For affine hyperbolic $S$, fill one new distinguished puncture
by $w:S_p\to S$.  The topological Birman sequence gives
$1\to\pi_1(S,p)\to M_{S_p}\to M_S\to1$.
Its completed kernel is exactly $P_S$: the composite
\[
 P_S\longrightarrow\widehat M_{S_p}
             \longrightarrow\operatorname{Aut}(P_S)
\]
is the inner-automorphism map and is injective, because a
nonabelian free profinite group is center-free.  Right
exactness of completion then gives
\[
 1\longrightarrow P_S\longrightarrow\widehat M_{S_p}
                \longrightarrow\widehat M_S\longrightarrow1. \tag{E.20}
\]
It is the right stabilizer of $w$, and therefore is preserved
by $F$.  This recovers an outer automorphism $\psi_S$ of $P_S$.
Two-deletion squares show that $\psi_S$ preserves the kernel
$N_c=\overline{\langle\!\langle I_c\rangle\!\rangle}$ of
filling any old labelled puncture $c$ whenever the resulting
surface is still affine hyperbolic.

For a degree-$d$ cover $f:S\to T$, remove an ordinary target
point $q$ and all $d$ points $D=f^{-1}(q)$, distinguishing
$p\in D$.  Write $f':S_D\to T_q$, $u:S_D\to S$ and
$v:T_q\to T$ for the resulting square.  An image of this
square under $F$ is still the correct puncture pullback.
Indeed, for the images $J'\subset P_{T_q}$ and $J\subset P_T$
of its two covering homomorphisms, commutativity up to a
chosen target conjugator gives
\[
 J'\subset v_*^{-1}(J),\qquad
 [P_{T_q}:J']=[P_{T_q}:v_*^{-1}(J)]=d.
\]
The subgroups are equal.  Its $d$ new source cusps therefore
lie over $q$ with degree one; no old target cusp is killed by
$v_*$.  This also identifies which cusps the square deletes.

We need one cancellation calculation.  If $u$ fills extra
ordinary source punctures and $q_u:M_{S_D}\to M_S$ forgets
them, then
\[
 \operatorname{Stab}_{M_{S_D}}(fu)
                  =q_u^{-1}\operatorname{Deck}(f).        \tag{E.21}
\]
For if $fu a=fu$, surjectivity on surface groups implies that
$q_u(a)$ induces the same outer action as an element of
$N_{\pi_1T}(f_*\pi_1S)/f_*\pi_1S$, the deck group.
The injective discrete Dehn--Nielsen--Baer representation
then gives (E.21); the reverse inclusion is immediate.
The finite deck group embeds in $\widehat M_S$.  One direct
verification uses a graph spine of its affine quotient:
\[
 H_1(S,\mathbb Q)\simeq\mathbb Q\oplus
     \mathbb Q[\operatorname{Deck}(f)]^{,r-1},\qquad r\geq2.
\]
The action is faithful and hence is detected modulo some
integer on integral homology.  Since $q_u$ is surjective,
its completed kernel is the closure of its kernel, and the
preimage in (E.21) is a finite union of kernel cosets.  Thus
(E.21) remains exact after completion and in arbitrary
completed markings.

Now characterize inside $P_T\subset\widehat M_{T_q}$ the
elements $b$ admitting a lift $a$ with
\[
                  b f'=f'a,\qquad ua=u.                  \tag{E.22}
\]
The chosen sheet makes the lift unique, because a deck
transformation fixing a sheet is the identity.  Discretely
(E.22) says exactly that the loop $b$ returns that sheet;
its subgroup is $f_*\pi_1(S,p)$.  The same holds after
completion.  The target lifting subgroup is open; on its
intersection with $P_T$, (E.21) says the forgotten lift lies
in the finite deck group.  The second equation in (E.22)
is the kernel of that finite continuous quotient.  Its
discrete subgroup is dense, proving that its closure is
$f_*P_S$.  Filling $D\setminus\{p\}$ identifies its lift
with point pushing in $S_p$.
All of (E.22), its chosen sheet and its fillings are retained
categorical data.  Applying $F$ gives
\[
                  (Ff)_*\psi_S=\psi_T f_*                \tag{E.23}
\]
up to one target inner automorphism.  This proves the needed
naturality rather than imposing it.

\emph{Genus-zero gauges and labelled ramification.}
For a genus-zero surface with $b\geq5$ punctures, every
coloured mapping class completion is center-free.  The pure
case follows from (E.20), starting with the free four-point
group.  For the coloured case its permutation acts faithfully
on the rational abelianization of the pure group:
\[
 H_1(\Gamma_{0,b},\mathbb Q)=
 \mathbb Q^{\binom b2}/
 \left\langle\sum_{j\ne i}e_{ij}:1\leq i\leq b\right\rangle.
                                                               \tag{E.24}
\]
The arrangement presentation with three punctures fixed at
$0,1,\infty$ proves this formula by its independent meridians.
Distinct $e_{ij}$ stay distinct for $b\geq5$: a difference of
two would have coefficients $a_i+a_j$ supported on two edges;
the complement is connected with an odd cycle, forcing all
$a_i=0$.  A permutation acting trivially fixes every pair and
is the identity.  The same argument over $\mathbb Q_\ell$
detects the finite permutation in the profinite group.

The hypothesis on $F$ consequently allows gauges making the
actions on $T,T_q$ and the filling $v$ all identities.
Explicitly, first fix the target action and then $v$ by its
single right orbit.  The remaining inner source conjugator
maps to the center of the target, hence lies in $P_T$;
gauging by its inverse fixes $v$ and the source action.
Thus $\psi_T$ in (E.23) may be taken to be the identity.
It follows that $Ff$ and $f$ have conjugate image subgroups,
so the same unlabelled cover orbit.
For a retained source cusp $c$, preservation of $N_c$ in
(E.23) also gives equality of the target normal closures of
$f_*I_c$ and $(Ff)_*I_c$.  In
$P_T^{\rm ab}=\widehat{\mathbb Z}^{b}/
\widehat{\mathbb Z}(1,\ldots,1)$ these have images
$e\widehat{\mathbb Z}[c_i]$ and
$e'\widehat{\mathbb Z}[c_j]$.
Distinct cusp lines are distinct over every $\mathbb Q_\ell$;
equality forces $i=j$ and then equality of all prime valuations
of the positive integers $e,e'$, so $e=e'$.
In particular a retained simple ramification point over a
distinguished branch value is recovered uniquely.  Passing
through the specified finite labelling cover gives the same
assertion with retained colours.

\emph{A Hurwitz cover surjecting on the higher-genus group.}
Fix $g\geq1$ and a stable marked type $(g,n)$.  Choose
$d\geq2g+2n+5$ and set $b=2d+2g-2$.
Let $\mathcal H$ parametrize degree-$d$ simply branched maps
to $\mathbb P^1$ with $n$ labelled simple ramification points,
their separately labelled values, and the remaining branch
values unordered; divide target coordinates by
$\operatorname{PGL}_2$.
Riemann's covering classification gives a finite etale branch
map $\mathcal H\to\mathcal M_{0,n,[b-n]}$: there are finitely
many transposition monodromy tuples, and distinct moving
branch values uniquely transport each cover.  Connected
transposition monodromy is $S_d$, so the deck group is trivial
for $d\geq3$.

We verify the stronger assertion that
$\mathcal H\to\mathcal M_{g,n}$ is smooth, onto and has
geometrically connected fibers.  For a fixed $(C;P_i)$, let
$D$ vary among reduced degree-$d$ divisors disjoint from the
$P_i$.  This is a smooth irreducible open subset of
$\operatorname{Sym}^d C$.  The spaces $H^0(C,\mathcal O(D))$
form a vector bundle; the $n$ independent conditions
$df(P_i)=0$ cut out a vector subbundle.
The conditions of exact simple poles at $D$, simple critical
points and distinct critical values are nonempty open
conditions.  Indeed Riemann--Roch gives surjection onto all
jets of total length at most $2n+4$ under the bound on $d$.
A multiple unspecified critical point imposes two independent
conditions over a one-dimensional locus.  Two unspecified
critical points with equal values impose three conditions
over a two-dimensional locus.  A coincidence involving one
prescribed critical point imposes two over a one-dimensional
locus; two prescribed values impose one proper hyperplane.
Zero pole coefficients and zero second derivatives at the
$P_i$ each impose proper hyperplanes.  These finitely many
incidence types cannot cover the vector space.
Every resulting function has unique pole divisor $D$.
This is precisely the chart of maps with infinity unramified,
and its fibers are irreducible.  The same bundles and jet
maps work in families, proving smoothness and surjectivity.
Every target admits such a coordinate, whose choices form a
connected open subset of $\operatorname{PGL}_2$; quotienting
coordinates retains geometric connectedness of the fibers.
The argument descends from finite labelling covers for colours.

It follows that $\mathcal H$ is connected and
\[
             \pi_1^{\rm et}(\mathcal H)\twoheadrightarrow
                          \widehat M_{g,n}.               \tag{E.25}
\]
Indeed a connected finite etale cover of the base stays
connected: a clopen decomposition of its pullback would,
by connected fibers and the openness of the smooth map,
give a clopen decomposition of the base cover.

Choose a cover in $\mathcal H$, remove all its branch values
and their preimages, and retain the $n$ ramification labels.
This gives $f:S\to T$ and a filling $u:S\to(C\setminus\{P_i\})$.
Here $b\geq5$ and $S$ has $b(d-1)$ punctures, so every
single deletion used above remains affine hyperbolic.
The genus-zero gauges and ramification calculation show
$Ff=f$ after a coloured source gauge.  The open target
lifting subgroup $H_f\subset\widehat M_T$ is
$\pi_1^{\rm et}(\mathcal H)$.  Its lift
$\ell_f:H_f\to\widehat M_S$ is unique, since the deck group
is trivial, and is characterized by $hf=f\ell_f(h)$.
Applying $F$ gives $F_S\ell_f=\ell_f$.
Gauge the filled object so that $F(u)=u$; this uses the single
right orbit of the filling and its surjective forgetting map.
The resulting action on $\widehat M_{g,n}$ fixes the image
of $H_f$, which is all of that group by (E.25).
Its outer class is therefore trivial.  This proves every
stable type in positive genus, including proper curves, and
the genus-zero types are the hypothesis.
\end{proof}

\subsection{An audit of the 2026 procongruence argument}

\begin{result}{E.14}{Refuted}{Three exact assertions used in the cited proof}
The following assertions in the opened 2026 sources are false
as printed.
\begin{enumerate}
\item In the proof of \cite[Proposition 6.9, p.~36]{ET_Boggi2026},
the full closed hyperelliptic mapping class group $\Upsilon_g$
is identified with $B_{2g+1}/\langle z^2\rangle$.
\item \cite[Lemma 6.16, pp.~42--44]{ET_Boggi2026} asserts
uniqueness of a genuine inertia-preserving extension from the
oriented multicurve stabilizer $H$ to $H_i$.
\item The proof of \cite[Lemma 2.23, p.~22]{ET_BF2026} identifies
\[
 \widehat P\Gamma_{1,2,\vec\gamma}\cap\widehat P\Gamma_{0,5}
          =\widehat P\Gamma_{0,5,q(\gamma)}.                \tag{E.26}
\]
Here the embedded five-point sphere group is the pure braid
subgroup of $\widehat B_4/Z$, and $q(\gamma)$ is the
Birman--Hilden image of a nonseparating curve.
\end{enumerate}
These counterexamples refute the specified proof assertions;
they do not refute the main automorphism-group theorems.
\end{result}
\begin{proof}
For the first assertion take $g=2$.  The braid quotient in
question has abelianization $\mathbb Z/40$, since $B_5$ has
abelianization $\mathbb Z$ and its full twist has exponent
$5\cdot4=20$.  The full genus-two group has abelianization
$\mathbb Z/10$.  For a direct presentation check use
\cite[Example 4.1, p.~194]{ET_CW2007}: with $m=2g+2$,
\[
 W=\sigma _1\cdots\sigma _{m-1}\sigma _{m-1}\cdots\sigma _1,
 \quad
 \Upsilon_g=B_m/\langle\!\langle W^2,\omega_m,
                                      [W,\sigma_i]\rangle\!\rangle.
                                                               \tag{E.27}
\]
All braid generators have equal abelian image.  The two
noncommutator relations have exponents $4(m-1)$ and
$m(m-1)$, giving order $(m-1)\gcd(4,m)$, hence $10$ at
$g=2$.  The proposed braid group instead describes the
Weierstrass-point stabilizer, of index $2g+2$, after capping
the chosen boundary.  Its map to the full group is inclusion
of that stabilizer, so the discrepancy is visible in the
actual geometric map as well as in abelianization.

For the second assertion use exactly the notation of Lemma
6.16: $\sigma=\{\gamma_1,\ldots,\gamma_{g-1}\}$,
$\sigma_i=\sigma\setminus\{\gamma_i\}$,
$H=\widehat P\Gamma(S,\partial S)_{\vec\sigma}$ and
$H_i=\widehat P\Gamma(S,\partial S)_{\vec\sigma_i}$.
The identity on $H$ has the extensions
$1$ and $\operatorname{Inn}(\tau_{\gamma_i})$ on $H_i$.
The twist $\tau_{\gamma_i}$ centralizes $H$, whereas it does
not centralize $\tau_{\beta_i}$ for a curve $\beta_i$ meeting
$\gamma_i$ once.  The two twists have on their handle the
matrices
\[
 A=\begin{pmatrix}1&1\\0&1\end{pmatrix},\qquad
 B=\begin{pmatrix}1&0\\-1&1\end{pmatrix};
 \quad ABA^{-1}\ne B\pmod3.
\]
Both extensions preserve all twist inertia.  Thus uniqueness
fails even for the identity parameter and already in genus
three with one boundary component.  More generally
$\operatorname{Inn}(\tau_{\gamma_i}^t)$, $t\in\widehat{\mathbb Z}$,
gives the same explicit ambiguity.  The ambiguities so obtained
for different $i$ are restrictions of the single global
inner automorphism by $\prod_i\tau_{\gamma_i}^{t_i}$:
the other twists are central in $H_i$.
This proves why genuine uniqueness is unnecessary for these
particular ambiguities, without claiming that they exhaust
the kernel of restriction.

For the third assertion put $a=\sigma _1$, $b=\sigma _2$,
$d=\sigma _3$ in $G_2=\widehat B_4/Z$ and
$v=ba^2b$ as in E.11.  Its strand permutation is the identity,
so $v\in\widehat P B_4/Z=\widehat P\Gamma_{0,5}$.
The braid relation gives $av=va$, so it stabilizes
$q(\gamma)$, whose twist is $a^2$.
On the homology of the filled torus, however,
\[
 q(v)=BA^2B=\begin{pmatrix}-1&2\\0&-1\end{pmatrix}
                         =-A^{-2}.
\]
It sends the primitive homology vector of $\gamma$ to its
negative.  Hence it reverses the orientation of $\gamma$ and
does not lie in the left side of (E.26).  The failure is
already detected modulo $3$.  The correct left side is the
kernel of the branch character on the right side; it has
index two, as shown by this very element.  Equivalently,
in the centralizer quotient of E.11, it is the intersection
of the two explicitly different characters
$v\mapsto0,d\mapsto1$ and $v\mapsto1,d\mapsto0$.
All the abelianization orders, braid commutators and matrices
in this proof are checked by \texttt{E\_braid\_models.py}.
\end{proof}

\begin{result}{E.15}{Proved}{Repair of the closed hyperelliptic lifting step}
For every $g\geq2$ there is a faithful continuous homomorphism
\[
 \widehat{\operatorname{GT}}\longrightarrow
                         \operatorname{Out}(\widehat\Upsilon_g)
                                                               \tag{E.28}
\]
lifting the standard genus-zero action on
$\widehat\Upsilon_g/\langle\iota\rangle
=\widehat\Gamma_{0,[2g+2]}$ and preserving hyperelliptic
nonseparating and separating twist classes with multiplier
$\lambda$.  Thus the existence portion of Proposition 6.9
can be repaired without its false point-stabilizer identification.
\end{result}
\begin{proof}
Use (E.27) with $m=2g+2$.  Let $R$ be the braid automorphism
$\sigma_i\mapsto\sigma_{m-i}$; it is conjugation by the
standard half twist.  With
$y_m=\sigma_{m-1}\cdots\sigma_1^2\cdots\sigma_{m-1}$,
we have $R(y_m)=W$ and $R(\omega_m)=\omega_m$.
The genuine standard action $\Phi_F$ sends $y_m$ and
$\omega_m$ to their $\lambda$th powers
\cite[Appendix, pp.~34--35]{ET_HS2000}.
Therefore $R\Phi_FR^{-1}$ sends $W$ and $\omega_m$ to
$W^\lambda$ and $\omega_m^\lambda$.
It preserves the normal subgroup in (E.27): $W^2$ and
$\omega_m$ go into it, while $[W,\sigma_i]$ maps to the
identity in the quotient because $W$ is central there.
Applying the same argument to the inverse gives equality
of the closed normal subgroups.  We obtain genuine
automorphisms of $\widehat\Upsilon_g$, multiplicatively in
$F$, and hence (E.28).

Modulo $W=\iota$ this is the standard sphere action conjugated
by the fixed inner automorphism $R$, so it has its standard
outer parameter.  Faithfulness follows from the sphere
action for $m\geq6$ and
\cite[Main Theorem]{ET_HS2000}, or
\cite[Corollary 2.8(i)]{ET_HMM2022}.
Every symmetric nonseparating twist is a conjugate of the
lift of an adjacent half twist.  A symmetric separating
twist cutting off genus $h$ is a conjugate of
$\omega_{2h+1}^2$ in the branched-cover presentation.
The standard braid action sends the corresponding full
twists to their $\lambda$th powers, and conjugation by $R$
preserves these conjugacy assertions.  This proves the
inertia statement, including the topological type of the
separating curve.
\end{proof}

\subsection{The genus-zero system and the manuscript's two fundamental groups}

\begin{result}{E.16}{Known}{The decorated genus-zero completion has automorphism group GT}
For $n\geq3$ retain the labelled completed exit categories, their open
objects, their boundary branches and procyclic twist subgroups in the
open groups, all relabellings, all forgetting maps, and all genus-zero
clutching maps, with their actual composition identifications.  Take
continuous equivalences preserving this data and commuting with
relabellings, modulo coherent natural isomorphism.  Their group is
canonically $\widehat{\operatorname{GT}}$ in the standard tangential
normalization.  The compatible action on all labelled strata and on
every twist has one common multiplier
$\lambda\in\widehat{\mathbb Z}^{\times}$; these are consequences of
that description.  Merely preserving the three peripheral subgroups
of the isolated four-point open group and commuting with its
permutation action instead gives
$\widehat{\operatorname{GT}}_0$.
\end{result}
\begin{proof}
Here the decorations name subgroup embeddings, not just abstract
isomorphism types of groups.  Write $G_n=\widehat P\Gamma_{0,n}$.
Harbater--Schneps, Main Theorem (a),(b), identifies the centralizer of
the permutation action in the inertia-preserving outer automorphism
group of $G_n$ with $\widehat{\operatorname{GT}}_0$ for $n=4$, and
with $\widehat{\operatorname{GT}}$ for $n\geq5$
\cite[pp.~1--3, Main Theorem]{ET_HS2000}.  Its parameter is the
actual pair $(\lambda,f)$ in the genuine braid formula
\[
 \sigma_1\longmapsto\sigma_1^\lambda,\qquad
 \sigma_i\longmapsto
 f(y_i,\sigma_i^2)^{-1}\sigma_i^\lambda f(y_i,\sigma_i^2).
                                                               \tag{E.29}
\]
Restriction along a forgetting map retains that pair; at four
points it is the inclusion
$\widehat{\operatorname{GT}}\hookrightarrow
\widehat{\operatorname{GT}}_0$.  Thus a compatible family on the
open groups has precisely one parameter, detected at five points.
The three-point group is trivial.  In particular a family cannot
acquire an independent four-point parameter.

We explain why this is also the indicated exit-category statement.
For a multicurve $\sigma$, the arrows from its stratum to the open
object form the continuous left $G_n$-set
$G_n/\overline\Delta_\sigma$.  Its automorphism group is
$N_{G_n}(\overline\Delta_\sigma)/\overline\Delta_\sigma$;
the twist stabilizer theorem identifies this with the boundary
object's automorphism group.  For a nested pair, composition with
a chosen target-to-open arrow embeds its transporter quotient
into this $G_n$-set.  These assertions follow first on the
discrete transporter category, and then on the cofinal finite
transporter quotients: membership and equality are precisely the
finite quotient coset conditions.  Genus-zero stabilizer closure
and cofinality are the comparison used in E.5--E.6, with the same
cutting description for arbitrary labelled genus-zero flags.

Consequently an equivalence inducing the identity on $G_n$ can
be made to fix one smoothing arrow at each object by a unique
object automorphism.  After this modification it fixes every
arrow, by the preceding embedding.  The only remaining
modifications come from $Z(G_n)$; this is trivial for $n\geq4$,
and $G_3=1$.  This proves both injectivity of restriction and the
absence of an extra coherence parameter.

For existence, (E.29) preserves every packet full twist and its
power with multiplier $\lambda$.  In the parenthesized braid
groupoid, the action sends a crossing to its $\lambda$th power
and a change of parentheses to its product with the evaluated
word $f$.  The two hexagon relations say exactly that these
assignments respect a crossing past a packet; the pentagon
says exactly that the two composites of changes of parentheses
agree.  Hence the assignments define a functor on the presented
groupoid and commute with substitution of packets.  Capping
uses the central full-twist relation, which is preserved since
its image is its $\lambda$th power.  This gives the action on
the spherical groupoid, on its boundary transporters, and on
their twist quotients.  Erasing strands gives forgetting;
substitution of packets and capping gives genus-zero clutching.
The same presentation proves their composition identifications
are respected.  Formula (1), the string-doubling calculation of
\S0.3, and the Appendix of \cite{ET_HS2000} are the explicit
braid realizations of these maps.  Passing to finite quotients
and then their inverse limit gives continuous functors.  The
inverse parameter supplies the inverse equivalence.  This
proves existence and all the asserted compatibilities.
\end{proof}

\begin{result}{E.17}{Proved}{Cote 143 uses geometric base change, and individual faithfulness}
\label{sec:E-manuscript}
In cote 143, scans 29--33, the group denoted
$\widehat\Gamma_{g,n}=\pi_1(\overline M_{g,n})$ is the
geometric fundamental group of the \emph{open} moduli stack:
the bar in this passage denotes base change to
$\overline{\mathbb Q}$.  Its arithmetic action is
\[
 G_{\mathbb Q}\longrightarrow
             \operatorname{Out}(\widehat P\Gamma_{g,n}).        \tag{E.30}
\]
For every stable type of positive modular dimension this
homomorphism is injective, as is its procongruence quotient
action.  For $(0,3)$ it is the trivial homomorphism to $1$.
With the usual pure mapping class groups in the unstable
cases, it is trivial for $(0,0),(0,1),(0,2)$ and faithful
for $(1,0)$.

The surface group $\widehat\pi_{0,3}=\widehat F_2$ occurring
in scans 29--30 is instead the kernel of the next universal
curve sequence.  More precisely
\[
 \widehat\pi_{0,3}\xrightarrow{\sim}
 \widehat P\Gamma_{0,4},\qquad
 \widehat P\Gamma_{0,3}=1.                                    \tag{E.31}
\]
Thus the tripod automorphism assertion and the four-point
modular-group assertion are related by this specified
isomorphism; they are not two assertions about the trivial
three-point modular group.
\end{result}
\begin{proof}
The transcription at scan 29 defines $M_{g,n}$ using smooth
proper curves with distinct markings.  At scan 30 it puts
$\pi_1(\overline M_{g,n})$ on the left of the arithmetic
exact sequence with quotient $G_{\mathbb Q}$.  Scan 31 does
the same after the labelled-to-unlabelled quotient.  Finally
scan 33 identifies the kernel of the geometric universal
curve sequence with $\pi_1(\overline X\setminus S_n)$.
These three actual maps establish the stated interpretation
of the bar.  No choice between the doubtful basepoint glyphs
$\xi$ and $j$ changes them.  We have checked the supplied
page transcriptions; there is no corresponding image file
for these five pages in the supplied package, so this is
a mathematical determination of the notation, not a new
claim to have visually deciphered the handwriting.

For comparison with the compactification used elsewhere in
this report, the inclusion induces
\[
 \widehat P\Gamma_{g,n}\twoheadrightarrow
 \pi_1^{\rm et}(\overline{\mathcal M}_{g,n,\overline{\mathbb Q}})
 =\widehat P\Gamma_{g,n}/
       \overline{\langle\!\langle\text{boundary twists}
                                  \rangle\!\rangle}.           \tag{E.32}
\]
The equality follows by extending a finite cover across each
normal-crossing chart: extension is equivalent to trivial
local meridian action, and descent from charts imposes
exactly the displayed closed normal subgroup.  At $(0,4)$
the target is $\pi_1(\mathbb P^1_{\overline{\mathbb Q}})=1$.
Thus replacing the bar of cote 143 by stable compactification
would change (E.30) to a nonfaithful map already in this
smallest positive-dimensional case.

The individual procongruence faithfulness theorem is
Hoshi--Tamagawa in the precise formulation and proof of
\cite[Theorem 2.5, pp.~8--9]{ET_Iijima2012}.  Its hypothesis
is $3g-3+n>0$, and its conclusion over $\mathbb Q$ is
injectivity.  Here is the kernel argument, to specify its
use.  Put $\Pi=\widehat\pi_{g,n}$,
$\Sigma=\check P\Gamma_{g,n}\subseteq\operatorname{Out}(\Pi)$,
and let $A$ be the arithmetic universal-monodromy image.
The nodal-monodromy theorem gives
$A/\Sigma\simeq G_{\mathbb Q}$.  The kernel of its outer
action on $\Sigma$ is
\[
 C_A(\Sigma)/Z(\Sigma).                                       \tag{E.33}
\]
The centralizer theorem used in that proof identifies
$C_{\operatorname{Out}^C(\Pi)}(\Sigma)$ with the finite
automorphism group of the tautological curve: it is trivial
outside the exceptional types, and in positive genus the
possible involution is the geometric central involution.
An exceptional symmetry not in the pure geometric group
permutes the marked inertia classes; it is therefore not
in $A$, which fixes their labels.  Hence (E.33) is trivial.
For genus zero the proof uses the faithful tripod action
and the genus-zero congruence theorem.  These are exactly
the cases and exclusions checked in the cited proof; no
genus-at-least-three congruence theorem is used.  The
equivariant quotient
$\widehat P\Gamma_{g,n}\twoheadrightarrow\Sigma$ sends an
inner action to an inner action, so faithfulness on
$\Sigma$ proves faithfulness of (E.30).

For the four sphere types with $n\leq3$, every
orientation-preserving homeomorphism fixing the markings
is isotopic to the identity, so the group is $1$.
For a torus with no mark, evaluation at a point gives the
fibration of identity components of diffeomorphism groups.
The evaluation map from translations to the torus is a
homeomorphism.  Thus its point-pushing map on components
is zero, and forgetting the single mark induces the
isomorphism
$P\Gamma_{1,1}\xrightarrow{\sim}P\Gamma_{1,0}
=\operatorname{SL}_2(\mathbb Z)$.
The algebraic forgetting morphism from pointed elliptic
curves to genus-one curves induces this same geometric
fundamental-group map and commutes with arithmetic base
change.  Its completed isomorphism therefore identifies
the two outer Galois actions, proving the last case from
$(1,1)$.  Finally the universal curve over the unique
labelled tripod is $\mathbb P^1\setminus\{0,1,\infty\}$;
this proves (E.31), including its arithmetic compatibility.
\end{proof}


\begin{result}{E.18}{Refuted}{The discrete outer-automorphism assertion, in every type}
For the pure orientation-preserving discrete mapping class group
$P\Gamma_{g,n}$ of a closed orientable surface with $n$ labelled
punctures, the complete table is
\[
\begin{array}{c|c}
(g,n)&\operatorname{Out}(P\Gamma_{g,n})\\ \hline
(0,n),\ n\leq3&1\\
(0,4)&\operatorname{GL}_2(\mathbb Z)\\
(1,0),(1,1)&C_2\times C_2\\
(1,2)&C_2\\
(2,0)&C_2\times C_2\\
\text{all remaining types}&S_n\times C_2.
\end{array}                                                     \tag{E.34}
\]
Here $S_0=S_1=1$.  Thus the question
``$\operatorname{Autext}(\Gamma_{g,n})=1$?'' of cote 143,
scan 32, has an affirmative answer exactly for the sphere
with at most three labelled punctures.
\end{result}
\begin{proof}
We use the precise finite-index rigidity theorem, not an
unqualified statement that all automorphisms are geometric.
Ivanov, Theorem 8.5.A, pp.~103--104 of the 1998/1999 text
\cite{ET_Ivanov1999}, states that an isomorphism between
finite-index subgroups of the extended mapping class group
is conjugation by an extended mapping class, except for
spheres with at most four punctures, tori with at most two
punctures, and the explicitly stated central-character
correction for the closed genus-two surface.  We opened the
actual theorem and its proof.  The proof first recognizes
powers of twists, then their commutation graph, and uses
Theorem 8.1.A to obtain a common conjugator; comparing its
action on every twist forces the remaining factor into the
kernel of the curve-complex action.  That kernel is trivial
in the ordinary cases and is the hyperelliptic $C_2$ in the
closed genus-two case.

Apply this theorem to $P\Gamma_{g,n}$, a normal subgroup
of index $2n!$ in the extended group.  In the ordinary
cases its centralizer there is trivial: a centralizing
mapping class fixes every curve, since it commutes with
every positive twist, and the curve action is faithful.
An orientation-reversing class cannot centralize a twist,
since it changes its exponent to $-1$.  Hence
\[
 \operatorname{Out}(P\Gamma_{g,n})
   =\Gamma_{g,[n]}^{\pm}/P\Gamma_{g,n}
   \xrightarrow{\sim}S_n\times\{1,-1\},                       \tag{E.35}
\]
where the displayed map records the actual puncture
permutation and orientation.  It is surjective: small
supported interchanges realize puncture permutations,
and a reflection can be chosen to fix all the punctures.
Its kernel is the pure orientation-preserving group by
definition.  This proves the last line, with its specified
comparison map.

For $(2,0)$ the same theorem gives, in addition, the central
transvection
\[
 T(h)=h\,\iota^{\chi(h)},
 \qquad \chi:P\Gamma_2\twoheadrightarrow C_2,                  \tag{E.36}
\]
where $\chi$ sends every nonseparating twist to $1$.
The presentation (E.27) with $m=6$ gives
$P\Gamma_2^{\rm ab}=\mathbb Z/10$ and
$\chi(\iota)=0$, since the word $W$ has length $10$.
Thus $T^2=1$.  The rigidity theorem says these and the
orientation automorphisms exhaust the possibilities.
The transvection is not inner: its square on each twist
is the original squared twist, so a putative conjugator
would fix every curve and be central; yet $T$ is nontrivial.
It commutes with orientation reversal, since $\chi$ is
unchanged modulo two.  This gives the second $C_2$ in that
row.  The same classification is stated in
Ivanov--McCarthy, Corollary 2, p.~4
\cite{ET_IvanovMcCarthy1999}.

For $(1,2)$ use the explicit isomorphism
$P\Gamma_{1,2}=B_4/Z(B_4)$ of E.9--E.11.  On the sphere
with five punctures this is the subgroup fixing the fifth
puncture, with the other four allowed to permute.  It is
therefore the inverse image of $S_4$ under
$\Gamma_{0,[5]}\twoheadrightarrow S_5$.  Apply Theorem
8.5.A to this finite-index subgroup of the five-punctured
sphere group.  Its normalizer in the extended sphere group
is itself together with orientation reversal, because a
point stabilizer is self-normalizing in $S_5$.  Its
centralizer is trivial by the same twist-power argument.
The quotient normalizer calculation gives exactly $C_2$.
In particular puncture interchange on the elliptic curve
creates no further outer class: the involution
$x\mapsto p+q-x$ supplies a natural isomorphism from the
label-swapped two-pointed curve back to the original.

For $(1,0)$ and $(1,1)$ the group is
$L=\operatorname{SL}_2(\mathbb Z)=C_4*_{C_2}C_6$, say
\[
 L=\langle r,s\mid r^4=1,\ r^2=s^3\rangle,
 \qquad z=r^2=s^3=-I.
\]
The center is $\langle z\rangle$ and the quotient is
$C_2*C_3$.  An automorphism of this free product preserves
the conjugacy classes of its two maximal finite factors.
After an inner modification, their two fixed vertices in
the Bass--Serre tree must be adjacent: otherwise the
subgroup generated by their stabilizers has a fundamental
segment of length greater than one and cannot be the
whole free product.  It follows that its outer
 automorphism group is $C_2$, represented by inversion of
the $C_3$ generator.  The kernel on outer automorphisms
from $L$ consists exactly of
$h\mapsto h z^{\chi(h)}$, with
$\chi\in\operatorname{Hom}(L,C_2)$ and $\chi(z)=0$.
There is one nontrivial such character:
$\chi(r)=1,\chi(s)=0$.  The two commuting involutions
\[
 (r,s)\longmapsto(r^{-1},s),\qquad
 (r,s)\longmapsto(r^{-1},s^{-1})                              \tag{E.37}
\]
represent the central transvection and orientation reversal.
The former is non-inner, since it induces the identity on
the center-free quotient, whose centralizer in itself is
trivial; an inner lift with that property would be central
in $L$.  The latter has nontrivial outer image on $C_2*C_3$.
This proves $\operatorname{Out}(L)=C_2^2$.

Finally $P\Gamma_{0,4}=F_2$ and abelianization induces
$\operatorname{Out}(F_2)\xrightarrow{\sim}
\operatorname{GL}_2(\mathbb Z)$.  The elementary Nielsen
basis operations give surjectivity; Nielsen's two-generator
reduction says two basis substitutions with the same four
exponent sums differ by simultaneous conjugation, giving
injectivity.  These are the reduction and concluding
identification in \cite[\S\S III--IV, pp.~395--397]{ET_Nielsen1917}.
For $n\leq3$ the pure sphere group is trivial.

There is also a uniform direct witness against triviality
in every positive genus.  A reflection fixing the markings
sends a nonseparating positive twist $T_v$ to $T_v^{-1}$.
An inner mapping-class automorphism induces symplectic
conjugacy $T_v\mapsto T_{Av}$.  Equality
$T_{Av}=T_v^{-1}$ over $\mathbb Z$ would force the primitive
vectors $Av=\pm v$ by their rank-one images, which gives
$T_{Av}=T_v$, a contradiction.  For the sphere with at
least four punctures, forgetting all but four points sends
an appropriate boundary meridian to a primitive generator
of $H_1(F_2,\mathbb Z)$; reflection negates its class,
whereas inner automorphisms act trivially on homology.
The finite $\operatorname{SL}_2(\mathbb F_3)$ check in
\texttt{E\_real\_extension\_checks.py} verifies the smallest
positive-versus-inverse twist obstruction.
\end{proof}

\begin{result}{E.19}{Proved}{The exact involution and centralizer fibres of LM6 C.15}
Let $P^{\pm}=P\rtimes\langle\tau\rangle$ be a profinite
group with $\tau^2=1$, let $K$ be a closed normal
$\tau$-stable subgroup of $P$, and put $Q=P/K$.
Write $\alpha(k)=\tau k\tau^{-1}$.  The compact pointed set
\[
 H^1(C_2,K)=
 \{k\in K:k\alpha(k)=1\}/
       (k\sim b k\alpha(b)^{-1},\ b\in K)
\]
has an action of $C_Q(\bar\tau)$ given by
\[
 q(g)\cdot[k]=[g k\alpha(g)^{-1}],\qquad g\in P.               \tag{E.38}
\]
The $P$-conjugacy classes of involutions of $P^{\pm}$ whose
images are $Q$-conjugate to $\bar\tau$ are exactly
\[
 C_Q(\bar\tau)\backslash H^1(C_2,K).                           \tag{E.39}
\]
For each cocycle $k$ there is an exact sequence of profinite
groups
\[
 1\longrightarrow K^{k\tau}\longrightarrow C_P(k\tau)
 \longrightarrow\operatorname{Stab}_{C_Q(\bar\tau)}([k])
 \longrightarrow1.                                           \tag{E.40}
\]
In particular the quotient centralizer need not fix the
trivial class, and the entire set $H^1$ cannot generally
replace its orbit set in (E.39).
\end{result}
\begin{proof}
An element of the coset $K\tau$ squares to $1$ exactly
when its coefficient is a cocycle.  Conjugation by $b\in K$
gives the displayed equivalence relation.  For a lift $g$
of an element centralizing $\bar\tau$,
\[
 g(k\tau)g^{-1}=gk\alpha(g)^{-1}\tau.
\]
The coefficient lies in $K$ and is again a cocycle.
Replacing $g$ by $bg$ changes its class by the specified
$K$-conjugation.  Replacing $k$ by an equivalent cocycle
has the same effect.  Composition of conjugations proves
the action law.  The cocycle locus is closed in $K$;
its quotient by the compact group $K$ is compact Hausdorff,
and the action is continuous with closed stabilizers.

An involution whose quotient is merely conjugate to
$\bar\tau$ can first be conjugated by a lift in $P$ to
make its quotient equal to $\bar\tau$.  The remaining
conjugators have images in $C_Q(\bar\tau)$, proving
(E.39).  The kernel of the centralizer map is $K^{k\tau}$.
If a quotient centralizer element fixes $[k]$, choose its
lift $g$ and an element $b\in K$ such that
$g(k\tau)g^{-1}=b(k\tau)b^{-1}$.  Then $b^{-1}g$ is a
centralizing lift.  This proves (E.40), including
surjectivity onto the stated subgroup.  Compactness gives
the asserted profinite exact sequence.

For a complete finite example take
$P^{\pm}=D_8=\langle r,t\mid r^4=t^2=1,trt=r^{-1}\rangle$,
$P=\langle r\rangle$, $K=\langle r^2\rangle$, and
$\tau=t$.  The action on $K$ is trivial, so $H^1$ has the
two classes $[1],[r^2]$.  The nontrivial element $q(r)$ of
$Q=C_2=C_Q(\bar t)$ interchanges them, because
$r\alpha(r)^{-1}=r^2$.  Thus (E.39) has one element,
whereas $H^1$ has two; $C_P(t)=K$ has trivial image in
$Q$.  The exhaustive eight-element multiplication-table
calculation in \texttt{E\_real\_extension\_checks.py} checks
all cocycles, conjugacy classes, the action, and the
centralizer.  This proves the qualification already
present in the actual statement of LM6 C.15.
\end{proof}

\subsection{Real centralizers in genus three}
\begin{result}{E.20}{Proved}{The five real genus-three centralizers with an oval}
Let $\tau$ be an orientation-reversing involution of a closed
oriented surface of genus three, with nonempty fixed locus.  Set
\[
 P=\widehat\Gamma_3,\qquad Q=\check\Gamma_3,\qquad
 K=\ker(P\longrightarrow Q),\qquad C=C_{\Gamma_3}(\tau).
\]
The quotient map induces a topological isomorphism
\[
 \overline C^{\,P}\xrightarrow{\sim}
 C_Q(\bar\tau)=\overline C^{\,Q},\qquad
 C_P(\tau)=K^\tau\rtimes\overline C^{\,P}.                    \tag{E.41}
\]
Thus LM6 C.17 is valid.  Its proof uses the finite-subgroup
centralizer theorem only for the cyclic group of order two.
\end{result}
\begin{proof}
We first establish the pointed comparison used in recognizing
an oval, since this also removes a dependence on the problematic
multitwist assertions of the 2020 source.  For a closed surface
$S$ of genus at least two put $\Delta=\widehat{\pi_1(S)}$.
Let $G_i=\widehat\Gamma_{g,i}$ for $i=0,1,2$, let
$V=\ker(G_2\to G_0)$, and let $F=\ker(G_2\to G_1)$.
The full configuration and Birman sequences identify
\[
 V=\widehat{\pi_1(\operatorname{Conf}_2(S))},\qquad
 F=\widehat{\pi_1(S\setminus\{p\})},\qquad V/F=\Delta.
                                                               \tag{E.42}
\]
Here is a direct check of the injectivity in these completed
sequences.  Conjugation gives a map from the discrete total
group to the automorphisms of the completed fiber group.  The
restriction to the fiber is its inner action.  The completed
fiber is center-free: for a punctured surface it is a nonabelian
free profinite group, and for $V$ this follows from its
center-free kernel $F$ and center-free quotient $\Delta$.
Consequently the completed fiber maps injectively into the
completed total group.  Right exactness of completion now gives
(E.42).  This argument applies as well to the one-point sequence.

Hoshi--Mochizuki's Theorem 6.1 and its proof
\cite{ET_HMNodal2011} say that the projection
\[
 \operatorname{Out}^{\mathrm{FC}}(V)
       \longrightarrow\operatorname{Out}(\Delta)
\]
is injective.  The outer actions supplied by the configuration
fibration preserve its fiber subgroups and its cuspidal inertia,
so they belong to this specified domain.  Suppose that
$z_1\in G_1$ acts trivially in $\operatorname{Aut}(\Delta)$,
and lift it to $z_2\in G_2$.  Its image in $G_0$ acts trivially
in $\operatorname{Out}(\Delta)$, hence trivially in
$\operatorname{Out}(V)$.  Multiplication by an element $v\in V$
therefore makes $z_2$ centralize $V$.  If $v_1\in\Delta$ is
its image in $G_1$, the resulting image $v_1^{-1}z_1$ acts
trivially on $\Delta$.  The original $z_1$ did too, so
$\operatorname{Inn}(v_1)=1$ and $v_1=1$.  Centralization of
$F$ consequently proves that $z_1$ acts trivially in
$\operatorname{Out}(F)$.

Conversely, geometric monodromy on $F$ preserves the conjugacy
class of its actual peripheral generator $c$.  An outer
automorphism with this property can be represented by one fixing
$c$.  Two such representatives differ by conjugation by an
element of $C_F(c)=\overline{\langle c\rangle}$, which becomes
trivial after filling $c$.  This defines the homomorphism
\[
 \operatorname{Out}(F;[c])\longrightarrow\operatorname{Aut}(\Delta).
\]
For the mapping-class action it is the point-pushing
conjugation map: this is the discrete filling construction, and
its equality extends to the completions by continuity.  Thus
trivial outer action on $F$ implies trivial action on $\Delta$.
The two kernels just computed are equal.  Since the image in
$\operatorname{Aut}(\Delta)$ contains
$\operatorname{Inn}(\Delta)$, lifting an outer automorphism
and then multiplying by the required inner automorphism proves
\[
 \check\Gamma_{g,1}\xrightarrow{\sim}
 \operatorname{Aut}(\Delta)
       \mathbin{\times}_{\operatorname{Out}(\Delta)}
                         \check\Gamma_g.                       \tag{E.43}
\]
The actual pointed reflection extends this isomorphism to the
orientation extensions.  This derivation uses the opened
configuration-group theorem, not a centralizer theorem for a
mixed-sign multitwist.

Let $A=\{A_1,\ldots,A_r\}$ be the oval multicurve.  Its
components are essential and pairwise nonisotopic.  A disk
bounded by an oval would be exchanged with the other side and
would force a sphere.  An annulus between two ovals would be
exchanged with a second complementary annulus and would force
genus one.  These possibilities are excluded here.
We next prove that every $B\in C_Q(\bar\tau)$ preserves these
actual unoriented curve classes.  The real surface extension is
\[
 E_\tau=\operatorname{Aut}(\Delta)
       \mathbin{\times}_{\operatorname{Out}(\Delta)}
                         \langle\bar\tau\rangle.
\]
This is the completion of the topological real surface
extension, by the center-free extension computation of E.8.
The conjugacy theorem for order-two subgroups and the
centralizer assertion of Boggi--Zalesskii's Theorem A
\cite{ET_BZ2024}, applied to this virtually surface group,
identify its section classes with the ovals.  The centralizer
in $\Delta$ of a section belonging to $A_i$ is
$\overline{\langle a_i\rangle}$, the full procyclic completion
of its oval group.  Surface groups are good and of finite
cohomological type; adjoining the reflection preserves
finite virtual cohomological type and the required cohomology
comparison.  Thus this use of Theorem A has the prime $2$ and
satisfies its discrete-module hypotheses.

For clarity, recognizing a procyclic subgroup alone would
leave an arbitrary profinite unit on its generator.  The
following marked-cover argument removes that ambiguity.
Choose one real point $p_i$ on $A_i$ and its pointed
reflection.  By (E.43), a lift of $B$ can, after an inner
adjustment in $\Delta$, be chosen to conjugate it to a
pointed reflection on some $A_j$.  The same conjugacy now
holds in the punctured procongruence group, and hence in its
action on the punctured surface group.

Here is the finite incidence construction used by that action.
For a finite cover of the real surface invariant under the
reflection, delete the inverse images of the marked point.
The order-two sections of the resulting real extension are
the components of its real locus.  At each deleted real point,
the local extension is
$\widehat{\mathbb Z}\rtimes C_2$, with inversion action.
Its two conjugacy classes of sections are the two branches,
since their coefficients are the two classes in
$\widehat{\mathbb Z}/2\widehat{\mathbb Z}$.  Intrinsically the
local extension is the normalizer of the real cusp inertia
in the punctured real extension.  Its inertia-conjugacy
classes of sections are the flags, and their images among
global section classes specify the adjacent interval edges.
For covers pulled back from the filled surface the local
inertia maps have degree one, so the resulting graph maps
are genuine graph coverings with their deck actions.  Inclusion of
these local extensions specifies the two incidences of that
point with the adjacent real intervals.  A selected real oval
therefore gives a finite cycle graph, with these branches as
flags.  These constructions use sections, centralizers,
peripheral inclusions and the maps between finite covers;
they are consequently preserved by the punctured group
isomorphism.  The one-marked base oval is a cycle with one
vertex and one edge and with two distinct flags.  An
isomorphism of that graph has degree
$\epsilon_i\in\{1,-1\}$.

Let $N$ range over reflection-invariant normal finite-index
subgroups of the filled surface group.  These form a cofinal
system, by intersecting an arbitrary subgroup with its
reflection translate and its normal core.  The component
above $A_i$ has graph-cover group
$\langle a_i\rangle/(\langle a_i\rangle\cap N)$.
Lifting one circuit of the base graph in every such cover,
with its flags retained, shows
\[
 B(a_i)\equiv a_j^{\epsilon_i}
             \pmod{\overline{\langle a_j\rangle}\cap B(N)}.
\]
The sign is independent of $N$ because the graph maps lie
over the same one-edge graph.  Cofinality proves
$B([a_i])=[a_j^{\epsilon_i}]$.  Thus $B$ preserves $A$ as
an actual unoriented multicurve in the procongruence curve
complex.  This supplies the generator-level step required
by the stabilizer theorem.

Put $L=\operatorname{Stab}_{\Gamma_3}(A)$ and let $L_0$
fix each oval and its orientation.  Cutting, and replacing
the boundary circles by labelled punctures, gives
\[
 1\longrightarrow T=\mathbb Z^r\longrightarrow L_0
    \longrightarrow\prod_j\Gamma(S_j)^{\mathrm{pure}}
    \longrightarrow1.                                        \tag{E.44}
\]
The five possibilities are
\[
\begin{array}{c|c}
(r,a)&\coprod_j S_j\\\hline
(1,1)&S_{2,2}\\
(2,1)&S_{1,4}\\
(3,1)&S_{0,6}\\
(2,0)&S_{1,2}\amalg S_{1,2}\\
(4,0)&S_{0,4}\amalg S_{0,4}.
\end{array}                                                    \tag{E.45}
\]
Here $a=0$ means dividing.  Euler characteristic gives the
genera displayed: a nondividing cut has genus $3-r$ and
$2r$ punctures; each half of a dividing cut has genus
$(4-r)/2$ and $r$ punctures.

Boggi's 2014 Theorem 4.5, including its proof and Lemma 4.6
\cite{ET_Boggi2014}, identifies
$Q_A=\overline L^{\,Q}$ and gives
\[
 1\longrightarrow\widehat T\longrightarrow Q_{\vec A}
 \longrightarrow\prod_j\check\Gamma(S_j)^{\mathrm{pure}}
 \longrightarrow1.                                           \tag{E.46}
\]
The full lattice $\widehat T=\widehat{\mathbb Z}^r$ injects.
In that proof the Looijenga-level monodromy first separates
every finite quotient of the twist lattice.  For a single
nonseparating curve the remaining kernel calculation is
performed in the HNN presentation of the surface group:
a representative which is identity on the cut subgroup sends
the stable letter $t$ to $t c^u$, by the cyclic centralizer
of the boundary generator $c$.  It is precisely a profinite
twist.  The separating calculation uses the corresponding
amalgam and the same centralizer.  Iterating this calculation
gives (E.46).  Neither a simultaneous orientation choice over
the whole curve complex nor the later mixed-sign multitwist
assertion is used.

All components in (E.45) have punctures and genus at most
two, so the punctured genus-at-most-two congruence theorem
(D.19, using Boggi 2009 Theorem 3.5 only in its punctured range)
identifies the right side of (E.46) with the full completion
of the right side of (E.44).  Completing (E.44) is right
exact; its completed left map is injective since its composite
with (E.46) is the displayed injection of $\widehat T$.
The comparison therefore gives
$\widehat{L_0}\simeq Q_{\vec A}$.  The finite permutation
and orientation quotient is unchanged on passing to the
closure, and hence
\[
 \widehat L\xrightarrow{\sim}Q_A.                            \tag{E.47}
\]
The map $\widehat L\to P$ is injective as well, since its
composite with $P\to Q$ is (E.47).

We give the cohomological hypotheses needed for the final
centralizer calculation.  Pure genus-zero groups are good
by the ordered-configuration free-group fibrations: their
sections, characteristic finite quotients of the free kernels,
and Hochschild--Serre comparison establish exactness on
completion and goodness inductively.  The group
$\Gamma_{1,1}=\mathrm{SL}_2(\mathbb Z)$ is virtually free.
The closed genus-two group is hyperelliptic; a finite-index
subgroup missing its central involution embeds as a
finite-index subgroup of the six-point sphere group, giving
its goodness by Shapiro's lemma.  The exact completed Birman
sequences, verified as in (E.42), then give goodness for
these groups with the punctures in (E.45).  The same
arguments give finite-type resolutions and finite virtual
cohomological dimension.  The spectral-sequence comparison
for (E.44)--(E.46) gives these properties for $L_0$ and then
for its finite extensions $L$ and
$L^\pm=L\rtimes\langle\tau\rangle$.  For the discrete
$\mathbb F_2$-modules in Theorem A, finite-type resolutions
and filtered unions reduce the comparison to finite
coefficients.  Residual finiteness follows from the faithful
congruence realization (E.47).

Apply Boggi--Zalesskii's Theorem A to this particular good
group $L^\pm$ and its subgroup $\langle\tau\rangle$ of
order two.  It gives
\[
 C_{Q_A}(\bar\tau)=\overline{C_L(\tau)}^{\,Q}.
\]
The oval argument shows that the entire centralizer in $Q$
lies in $Q_A$; its discrete counterpart shows
$C_L(\tau)=C$.  This proves the first equality in (E.41).
The injection $\widehat L\to P$ and (E.47) identify the
two closures of $C$.  Finally every $u\in C_P(\tau)$ has
the same image as a unique $c\in\overline C^{\,P}$;
$uc^{-1}$ belongs to $K^\tau$.  Their intersection is
trivial, conjugation by the second factor preserves the
first, and compactness makes the resulting unique
factorization a topological semidirect product.  This
proves (E.41) in full.  E.19 then shows that the quotient
centralizer fixes the trivial $H^1$ class for these five
types.  The argument makes no assertion that $K^\tau$
vanishes or that the other $H^1$ orbits have been computed.
\end{proof}

\subsection{The intrinsic arithmetic condition}
\begin{result}{E.21}{Proved}{What the AA condition reconstructs from the system}
Let $X=\mathbb P^1\setminus\{0,1,\infty\}$ and let $X_2$
be its ordered two-point configuration space.  The AA
condition of Hoshi--Mochizuki--Tsujimura Definition 5.12 is
an intrinsic predicate on closed subgroups of the
$\widehat{\operatorname{GT}}$ reconstructed from
\[
 \Pi_{X_2}=\widehat\Gamma_{0,5}
       \longrightarrow\Pi_X=\widehat\Gamma_{0,4}.             \tag{E.48}
\]
Neither a selected arithmetic fundamental-group extension
nor an identification of a reconstructed field with a
preselected algebraic closure is an input to this predicate.
Its maximal closed subgroups give the conjugacy class of
$G_{\mathbb Q}$; adjoining a specified complex-conjugation
involution selects exactly the subgroup in E.7.
This answers the margin question in cote 143, scan 30,
using the configuration structure available at the tripod.
\end{result}
\begin{proof}
We record the data in the definition in order to specify
precisely what the reconstruction uses.  Corollary 4.5 and
Remark 4.5.1 of \cite{ET_HMT2025}, together with
Definitions 3.3, 4.1 and 5.1, construct from $\Pi_{X_2}$:
the generalized fiber quotients, their cuspidal inertia,
$\widehat{\operatorname{GT}}$, and, for a closed subgroup
$B$, its cofiltered arithmetic Belyi diagrams.  These are
diagrams of profinite groups over open subgroups of $B$;
``arithmetic'' in their name does not supply an arithmetic
extension as extra input.  Their cusps give a set $Q_B$
with constructed addition and multiplication.  ArBC requires
these operations to make $Q_B$ an algebraically closed field
of characteristic zero algebraic over its prime field.
This expresses the isomorphism type of
$\overline{\mathbb Q}$ without choosing an isomorphism to it.

The four additional clauses of Definition 5.12 use the
following constructed objects: the Kummer map for $Q_B$;
the inverse limit $\Pi$ of the Belyi-diagram groups; its
open subgroups $\Pi^\dagger$ and finite cusp fillings;
their degree-one cohomology and cuspidal divisor maps; and
sets $L_{\Pi^\dagger}$ of functions on the reconstructed
cusps.  They require an infinite divisible subfield,
the specified cohomological vanishing and divisor
injectivity, compatible subfields
$K_{\Pi^\dagger}\subset L_{\Pi^\dagger}$ with
$K_\Pi=Q_B(t)$ and the prescribed finite Galois groups
and Kummer classes, and a lift of the reconstructed
tripod isomorphism to its two-point configuration group.
The last clause is the commutative diagram
\[
\begin{array}{ccc}
 \Pi_{X_2}&\xrightarrow{\sim}&\Pi_{(X_{Q_B})_2}\\
 \downarrow&&\downarrow\\
 \Pi_X&\xrightarrow{\sim}&\Pi_{X_{Q_B}},
 \qquad X_{Q_B}=\operatorname{Spec}Q_B[t,t^{-1},(t-1)^{-1}].
\end{array}                                                    \tag{E.49}
\]
Remark 5.12.2 explicitly treats the left column as an
abstract topological group with its specified quotient.
The field and scheme on the right are outputs of the
preceding constructions.  Isomorphic input diagrams
transport all open subgroups, conjugacy classes, cohomology
maps, functions and field structures in these clauses;
they therefore preserve the AA predicate.  This is the
claimed intrinsic character of the construction.

The isomorphisms in (E.48) require no choice beyond the
labelled tower: fixing its first three marks at
$0,1,\infty$ identifies $M_{0,5}$ with $X_2$, and
forgetting the fifth mark is its first projection to
$M_{0,4}=X$.  The arithmetic Belyi diagrams just used
range over open subgroups and their morphisms, all of
which are available from the profinite groups; adding
such diagrams does not assume an existing $G_{\mathbb Q}$
action.  If one starts with the abstract group
$\Pi_{X_2}$ alone, Corollary 4.5 reconstructs the
relevant quotients up to their stated finite relabelling
ambiguity.  The marked tower supplies the chosen quotient.

For comparison with the manuscript's type labels, its
surface group at $(0,3)$ is
$\Pi_X\simeq\widehat\Gamma_{0,4}$, by E.17; its modular
group $\widehat\Gamma_{0,3}$ is trivial.  Equation (E.48)
is the exact additional map used in AA.  Thus the actual
definition uses the two-point configuration data, and
makes no claim that the trivial modular group at $(0,3)$
alone determines the arithmetic subgroup.  Nor does
Definition 5.12 replace (E.49) by an unproved condition
on the three peripheral subgroups of $\Pi_X$ alone.

Theorem 5.17(ii), whose proof passes through Theorem
5.15(iv) and the compact-group inclusion lemma 5.16,
identifies the maximal AA subgroups with the
$\widehat{\operatorname{GT}}$-conjugates of
$G_{\mathbb Q}$.  This is an intrinsic conjugacy class
before a real framing is chosen.  For a fixed complex
conjugation $c$, the conjugacy of involutions inside
$G_{\mathbb Q}$ and
$C_{\widehat{\operatorname{GT}}}(c)=\langle c\rangle$
give uniqueness exactly as proved in E.7.  The chosen
$c$ is the additional framing that selects one member
of that class.  AA is a predicate on closed subgroups;
the resulting membership test for an element is
membership in this selected maximal subgroup.  It is
not an unprovided elementwise version of the four
clauses above.
\end{proof}

\subsection{The first unresolved elliptic compatibility}
\begin{result}{E.22}{Proved}{An exact word for the nonseparating elliptic comparison}
In the notation of E.11 set
\[
 R=\langle v,d\mid(vd)^2=1\rangle^{\wedge}
       \simeq C_2\amalg\widehat{\mathbb Z},\quad
 h=f(v,d^2),\quad k=f(d,v^2).
\]
Let $\delta_2:\ker(\widehat F(x,y)\to C_2)\to
\widehat F(x,y)$ send $x^2\mapsto x$, $y\mapsto y$,
and $xyx^{-1}\mapsto1$, where the character sends
$x\mapsto1,y\mapsto0$.  Write $c_2(f)$ for the exponent
of $y$ in $\delta_2(f)$.  The action induced by
$F=(\lambda,f)\in\widehat{\operatorname{GT}}$ on the
pure normalization group of the nonseparating boundary
of $M_{1,2}$ has the same standard genus-zero outer
parameter $F$ exactly when
\[
 f(d,v^2)d^{-4c_2(f)}f(v,d^2)=v^{4c_2(f)}\quad\hbox{in }R.
                                                               \tag{E.50}
\]
Every inner conjugacy implementing this comparison is the
one represented by $d^{-4c_2(f)}h$.  Equation (E.50) is
equivalent to the following equality in the full profinite
three-strand braid group, with standard generators $u,w$:
\[
 f(u,w^4)=w^{8c_2(f)}f(u^2,w^2)u^{4c_2(f)}
                           (uw)^{-6c_2(f)}.                   \tag{E.51}
\]
\end{result}
\begin{proof}
The standard pure subgroup of $R$ is
$K_{\rm std}=\langle v,d^2\rangle$, whereas the geometric
branch-labelled subgroup is
$K_{\rm geo}=\langle d,v^2\rangle$; these are the two
index-two kernels computed in E.11.  Each is free
profinite of rank two, as follows either from
Reidemeister--Schreier rewriting or its three-punctured
sphere cover.  The automorphism $S$ of $R$ interchanging
$v$ and $d$ interchanges these two subgroups.
The genuine action from E.11 is
\[
 \phi(v)=v^\lambda,\qquad \phi(d)=h^{-1}d^\lambda h.
\]
The standard action on the ordered generators $(d,v^2)$
of $K_{\rm geo}$ is the restriction of
$\psi=S\phi S$, with
$\psi(d)=d^\lambda$ and $\psi(v)=k^{-1}v^\lambda k$.
Thus the asserted parameter comparison is precisely
inner equivalence of these restrictions by an element
of $K_{\rm geo}$.

The free-product centralizer theorem gives
$C_R(d)=\overline{\langle d\rangle}$ and
$C_R(v)=\overline{\langle v\rangle}$; the second also
follows by applying $S$ to the first.  Since $\lambda$
is a unit, the same holds for the indicated
$\lambda$-powers.  The subgroup $K_{\rm geo}$ has
trivial centralizer in $R$.  For example a centralizer
element centralizes both $d$ and $v^2$; their maximal
procyclic centralizers have trivial intersection in
the free product.  Consequently two automorphisms of
$R$ agreeing on $K_{\rm geo}$ agree everywhere:
for each $r$, their images of $r$ differ by an element
centralizing that normal subgroup.

The equation $\psi=\operatorname{Inn}(g)\phi$ on $d$
forces $g=d^t h$ for $t\in\widehat{\mathbb Z}$.
Its equation on $v$ then forces
\[
                k d^t h=v^s.
\]
Since $h,k\in R'$, abelianization in
$R^{\rm ab}=\widehat{\mathbb Z}[d]\oplus C_2[vd]$
gives $s=-t$ and $t\in2\widehat{\mathbb Z}$.
Conversely such a solution implements the required
inner comparison, and $g$ belongs to $K_{\rm geo}$
as its character is zero.  Triviality of the center
of $R$ proves uniqueness of $g$ and hence of $t$.

We now determine $t$, rather than leave an unspecified
conjugating exponent.  Write $f$ in the Schreier basis
$(x^2,y,xyx^{-1})$ of the index-two subgroup of
$\widehat F(x,y)$.  Its exponent vector is
$(0,b,-b)$, where $b=c_2(f)$: the zero exponent of
$x$ in $f$ forces the first coordinate to vanish,
and the zero exponent of $y$ forces the other two
to sum to zero.  Put $X=v$, $Y=d^2$, the basis of
$K_{\rm std}$.  Under $x\mapsto d,y\mapsto v^2$
the three Schreier generators go to
\[
 Y,\qquad X^2,\qquad (X^{-1}Y^{-1})^2,
\]
because $dvd=v^{-1}$ in $R$.  Thus
$[k]_{K_{\rm std}}=(4b,2b)$, while
$[h]_{K_{\rm std}}=0$.  Abelianizing
$kd^th=v^{-t}$ in this free subgroup gives
$(4b,2b+t/2)=(-t,0)$, so $t=-4b$.
This proves (E.50), including both directions and
its exact conjugator.

The mixed subgroup $\langle u,w^2\rangle$ of
$\widehat B_3$ has central quotient $R$ under
$u\mapsto d,w^2\mapsto v$; its center is generated
by $(uw)^3=(uw^2)^2$.  Applying that quotient to
(E.51), and using $f(d^2,v)=h^{-1}$, gives (E.50).
Conversely (E.50) puts the discrepancy in (E.51)
in this central procyclic group.  Both sides of
(E.51) have exponent sum zero in $\widehat B_3$,
whereas its central generator has exponent sum six.
Multiplication by six is injective on
$\widehat{\mathbb Z}$, so the discrepancy is trivial.
This proves the equivalence with (E.51).
\end{proof}

\begin{result}{E.23}{Open}{The precise remaining whole-tower and elliptic equations}
The four requested rows have the following proved portions
and remaining scope:
\begin{center}
\begin{tabular}{p{0.23\linewidth}|p{0.68\linewidth}}
LM3b E.9 & Proved by E.8--E.10, with the specified
peripheral and point-pushing structure.\\
LM3b H.9 at $\overline M_{0,5}$ & Proved by E.5--E.6:
all flag stabilizers are good and have the required full topology.\\
LM3b E.4 & Orbifold normalizer comparisons and elliptic
forgetting are proved by E.11--E.12. The first nonseparating
clutching comparison for arbitrary GT is exactly (E.50).\\
LM3 F.16 and Lead 5(d) & The finite-sheaf completion
comparison holds in every stable type. A complete classification
of the automorphisms of the full decorated system remains
undetermined; in the cover-and-filling enhancement E.13 proves
injectivity of restriction on the induced outer-action families.
E.27 proves the strict marked-pants operad kernel is trivial,
and E.29 extracts its full completed cutting groups.
\end{tabular}
\end{center}
In particular no proof of the equality of the whole-system
automorphism group with $\widehat{\operatorname{GT}}$ is
claimed here.  The smallest missing existence comparison
in this approach is $\overline M_{1,2}$ along its
nonseparating boundary, with the two explicit maps of E.22.
\end{result}
\begin{proof}[Verified scope and exact unresolved steps]
Let $\mathcal T$ denote the full system of strict
completions, with the labels, strata, inertia embeddings
and all forgetting and clutching diagrams of E.16.
Restriction and the arithmetic construction give actual
homomorphisms
\[
 G_{\mathbb Q}\longrightarrow\operatorname{Aut}(\mathcal T)
       \xrightarrow{\mathrm{res}_0}\widehat{\operatorname{GT}},
                                                               \tag{E.52}
\]
whose composite is the usual faithful inclusion by E.4
and E.16.  The image therefore contains this specified
arithmetic subgroup.  E.13 proves that the kernel on
induced outer group actions is trivial after imposing
its complete covering-and-filling structure.  It does
not construct preimages of arbitrary elements of GT.
For the system without that enhancement, that argument
does not identify the kernel of $\mathrm{res}_0$.
These are the exact maps in the two versions of the
question.

E.27 proves a further unconditional kernel result
for the strict marked-pants surface operad:
$\operatorname{Aut}_0(\widehat{\mathbf S})\to
\widehat{\operatorname{GT}}$ is injective.
There is an explicit discrete extraction from the
boundary data: attach $n$ labelled tripods, each
with two distinct external labels, to a core of
type $(g,n)$, and take its cutting stabilizer.
The boundary-matching quotient of
$\Gamma_{g,n}^{\partial}\times
(\Gamma_{0,3}^{\partial})^n$, after killing the
$2n$ external boundary twists, is
$\Gamma_{g,n}^{\partial}$, since
$\Gamma_{0,3}^{\partial}=\mathbb Z^3$.
E.29 carries this extraction through full completion
by retaining both endpoints of the distinguished
crossing arrow.  Its diagonal stabilizer is the
full completed cutting group, even when its
projection to the open mapping class group has
unknown kernel.
To recover an automorphism of the exit tower
from its restrictions also requires recovery of
every crossing-coset map: for the left action on
$\widehat\Gamma/\overline\Delta_\sigma$ its
equivariant automorphism group is
$N_{\widehat\Gamma}(\overline\Delta_\sigma)/
\overline\Delta_\sigma$.
Thus the proposed recovery by stabilizer gauges
requires the exact equality
$N_{\widehat\Gamma}(\overline\Delta_\sigma)=
\overline{\operatorname{Stab}(\sigma)}$.
E.26 proves the two genus-one equalities used
in E.27.  No proof for every higher-genus
crossing subgroup or an explicit inverse
reconstruction of all its link maps is supplied
here; this is the remaining faithfulness step
for applying (E.66) to $\mathcal T$.
E.28 further makes the exact elliptic existence
comparison the fixed-point equation for its
explicit involution $\mathcal J$ of GT; the
identity $\mathcal J=1$ remains unproved.

The nonsolvable finite-subgroup centralizer equality
is not an input to (E.52), by E.1--E.4.  The full-to-
procongruence map is the quotient by the kernel of
$\widehat\Gamma_{g,n}\to\operatorname{Out}
(\widehat\pi_{g,n})$.  D.19 identifies it as an
isomorphism in genus at most two.  Thus it is already
an isomorphism at the smallest elliptic comparison in
E.22: that remaining equation cannot be removed merely
by replacing the full completion by the procongruence
completion.  E.19--E.20 specify the exact involution
fibers and centralizers of this quotient in genus
three, with no vanishing assertion about its kernel.

The opened proof of Boggi's 2026 claimed procongruence
automorphism theorem has the three precise defects
proved in E.14.  E.15 repairs its closed hyperelliptic
existence step.  The inner-twist counterexamples to
genuine extension uniqueness can themselves be
absorbed by a global inner conjugacy, as E.14 shows;
this does not prove that they exhaust the ambiguity.
Moreover the oriented-stabilizer error in
Boggi--Funar Lemma 2.23 lies upstream of the gluing
criterion used in that proof.  A full repair at those
steps has not been supplied.  Thus that proof does
not presently close the surjectivity or kernel
assertions in (E.52) in this audit.

The harmonic attack yields an exact further comparison.
Nakamura--Schneps relation (IV) \cite{ET_NS2000} is (E.51) with its
parameter $\rho_2(F)$.  Their Definition 5.1 defines
$\rho_2$ from the unique proword
$g(y,x)^{-1}g(x,y)=f(x,y)$, by
$g^{\rm ab}=(xy)^{\rho_2}$.  Their Proposition 5.3,
with its proof, identifies $\rho_2=c_2$ already on the subgroup
$I\Gamma\prime$ defined by (III').  To check the coefficient,
write $f=A_F(\bar x,\bar y)*[x,y]$ modulo the second
commutator subgroup.  The Schreier homomorphism $c_2$
takes $[x,y]$ to $-1$, changes sign under conjugation
by $x$, and is invariant under conjugation by $y$.
Hence $c_2(f)=-A_F(-1,1)$, exactly the coefficient
in their Proposition 5.3.  Imposing (IV) then gives
their subgroup $I\Gamma$.
Their construction therefore proves the required
clutching on that subgroup.  Arithmetic clutching is
also supplied directly by E.4.  The unresolved
assertion is (E.51) for all GT with the explicitly
computed coefficient $c_2(f)$.

A separate attempted derivation from Kummer distribution
has a verified endpoint.  Enriquez's Theorem 8.1 in
\cite{ET_Enriquez2007} proves, in the full profinite
topology,
\[
                  \delta_2(f)=y^{c_2(f)}f(x,y).             \tag{E.53}
\]
Its proof, including Lemmas 8.2--8.5, pulls the pentagon
to the configuration cover, applies its filling map,
and erases the zero strand.  The resulting equality
is exactly (E.53).  To finish the elliptic attack one
must derive the two-root substitution (E.51) from this
identity and the GT relations.  In Schreier coordinates
this requires controlling
$f(x,y^2)$ at
$(x^2,y^2,xy^2x^{-1})$, rather than only its quotient
that kills the third Schreier generator.  This is the
remaining algebraic step, not a change of completion.
The identity and complex conjugation satisfy it
($f=1,c_2=0$), as does every arithmetic element by
(E.52).  No genuine GT element violating it has been
constructed in this work.

The elliptic-associator route has also been checked
at the group level.  Put $c=[Y,X]$, $m=(\lambda-1)/2$.
Enriquez's Proposition 2.20 \cite{ET_Enriquez2014}
constructs the actual automorphism
\[
 A_F(X)=f(X,c)X^\lambda f(X,c)^{-1},\qquad
 A_F(Y)=c^m f(YX^{-1}Y^{-1},c)Y f(X,c)^{-1}.
                                                               \tag{E.54}
\]
For $T_+(X,Y)=(X,YX)$ and
$T_-(X,Y)=(XY^{-1},Y)$, let $\mathcal B$ be the
image of $\widehat B_3$ in
$\operatorname{Aut}(\widehat F(X,Y))$ generated by
these two automorphisms.  Conjugation in the conventions
of that paper is $T\mapsto A_F^{-1}TA_F$.
Its Remark 2.21 computes the images of $T_+$ and
$z=(T_+T_-)^3$, but normalization of $\mathcal B$
still requires the concrete generator test
\[
                   A_F^{-1}T_-A_F\in\mathcal B.             \tag{E.55}
\]
Proposition 8.6(b) proves the arithmetic case by
Nakamura's formula for the action on the universal
elliptic fiber; its proof uses that arithmetic
formula at its lower front square.  The section in
(E.54) and the fixed-elliptic-curve braid action alone
have not established (E.55) for all GT here.  An
identification of this route with the conjugator in
(E.50) also remains to be proved.  These explicit
maps record where this second approach stops.
\end{proof}

\begin{result}{E.24}{Open}{The arithmetic image, with exact indices and witness tests}
For the geometric modular groups in cote 143, let
$A_{g,n}$ be the subgroup of their outer automorphisms
preserving the boundary-inertia conjugacy classes and
centralizing the specified relabelling action.  In genus
zero the exact targets and arithmetic maps are
\[
\begin{array}{c|c|c}
 n& A_{0,n}&\rho_{0,n}:G_{\mathbb Q}\longrightarrow A_{0,n}\\\hline
 0,1,2,3&1&\text{trivial}\\
 4&\widehat{\operatorname{GT}}_0&\text{the tripod action}\\
 n\geq5&\widehat{\operatorname{GT}}&\text{the standard inclusion}.
\end{array}                                                     \tag{E.56}
\]
The surface-tripod formulation occurs at surface type
$(0,3)$ through
$\widehat\pi_{0,3}\simeq\widehat\Gamma_{0,4}$.
For $n\geq5$, surjectivity in (E.56) is exactly the
unresolved equality $G_{\mathbb Q}=\widehat{\operatorname{GT}}$.
The same equality describes the arithmetic image in
$\operatorname{Out}(\widehat\Gamma_{1,2})$, whose group
is GT and whose boundary inertia is preserved by E.11.
The remaining positive-genus targets are not classified
by the results of this section; E.17 establishes their
individual arithmetic faithfulness in its stated range.

For the element designated in LM7 E.58, the tests for
peripheral preservation, $\mathfrak S_3$-commutation,
multiplier one and trivial peripheral homology all pass.
Its specified modulo-twelve frame certificate has CM
parity one, whereas the arithmetic parity is zero.
It consequently fails membership in the selected maximal
AA subgroup of E.7 and E.21.  Its membership in GT,
and the value of the supplied global character $\Delta$
on all of GT, remain undecided here.
\end{result}
\begin{proof}[Proof of the target identifications and the verified tests]
The three-point labelled modular stack is a point;
normalizing its three markings also makes all lower
sphere modular groups trivial.  Four labelled points
have moduli $X=\mathbb P^1\setminus\{0,1,\infty\}$,
and the action of $\mathfrak S_4$ on its group factors
through $\mathfrak S_3$.  Harbater--Schneps' Main
Theorem (a) identifies its inertia-preserving
$\mathfrak S_3$-centralizing outer group with
$\widehat{\operatorname{GT}}_0$, including both
directions of the pair construction.  For $n\geq5$,
Hoshi--Minamide--Mochizuki Corollary 2.8 identifies
\[
 \operatorname{Out}(\widehat\Gamma_{0,n})
       =\widehat{\operatorname{GT}}\times\mathfrak S_n.
\]
The relabellings are the second factor, whose center
is trivial.  Its centralizer is therefore the first
factor; its inertia preservation is given by the
explicit braid action of E.16.  Forgetting all but
four marks gives the standard inclusion into GT$_0$.
These maps prove (E.56) and locate the first modular
GT surjectivity problem at $n=5$, with no change of
the manuscript's surface group into a modular group.
The elliptic assertion follows from the isomorphism
$\operatorname{Out}(\widehat\Gamma_{1,2})=\widehat{\operatorname{GT}}$
used and checked in E.9--E.12.  Relabelling its two
marks is inner, and E.11 covers its two boundary types.

Here is an independent finite verification of the
E.58 test.  Write
\[
 B=C_2\amalg C_3=\langle s,r\rangle,\qquad
 F_2=\ker(B\to\mathfrak S_3).
\]
\[
 U=\begin{pmatrix}1&1\\0&1\end{pmatrix},\quad
 L=\begin{pmatrix}1&0\\-1&1\end{pmatrix},\quad
 S=ULU,\quad R=-UL.
\]
The signed image of $F_2$ modulo twelve is generated
by $U^2,L^2$.  It has order $96$, has a four-element
image modulo four, and does not contain $-I$.
The canonical lift of the prescribed projective
frame $[5I]$ is therefore $5I$: it is $I$ modulo
four and $-I$ modulo three.  The designated witness
has normalized frames
\[
                    T=5I,\qquad W=I.                       \tag{E.57}
\]
In $\mathrm{SL}_2(\mathbb F_3)$, the centralizer
of $R$ is $\{\pm R^i:0\leq i<3\}$ and its
fixed line is $\mathbb F_3(1,-1)$; the centralizer
of $S$ is $\{\pm I,\pm S\}$.  Define
$\epsilon_3(T)$ as the sign on the first line,
and $\epsilon_2(W)$ as the quotient
$\langle S\rangle/\langle-I\rangle$.
Thus (E.57) gives
\[
             \epsilon_3(T)=1,\quad\epsilon_2(W)=0,
             \qquad\epsilon_3(T)+\epsilon_2(W)=1.          \tag{E.58}
\]
The independent program \texttt{E\_CM\_checks.py}
enumerates every matrix involved.  It finds frame
centralizers of orders $6$ and $8$, with exactly
$24$ of their $48$ pairs of each parity; it also
checks $[U^2,L^2]^2=5I$ modulo twelve.  This is a
finite calculation actually run, not an appeal to
the referee's scripts.

For the arithmetic comparison put
$E_3:y^2=x^3+1$, $E_4:y^2=x^3-x$ and
$r_0=(3/4)^{1/4}$.  The precise field identity is
\[
       \mathbb Q(E_3[4])=\mathbb Q(i,r_0)
                       =\mathbb Q(x(E_4[3])).             \tag{E.59}
\]
One can verify both inclusions directly.  The halves
of $(-1,0)$ on $E_3$ have $x=-1\pm\sqrt3$;
one $y$-coordinate is
$\sqrt3(\sqrt3-1)r_0$, and multiplication by
$\zeta_3$ gives the other nonzero two-torsion halves.
Conversely $\sqrt3$ is obtained from these
$x$-coordinates, $r_0$ from the displayed
$y$-coordinate, and $i$ from the Weil pairing.
For $E_4$, the three-division polynomial is
$3x^4-6x^2-1$.  Choose roots $a,b$ with
$a^2=1+2/\sqrt3$, $b^2=1-2/\sqrt3$ and
$ab=i/\sqrt3$.  Then $w=a+ib$ satisfies
$w^2=2/\sqrt3$, so $w=r_0^{-1}$ after a choice
of fourth root.  Also
$a=(w+2/w)/2$ and $b=(w-2/w)/(2i)$.
These formulas give both field inclusions in the
second equality; $\sqrt3$ and $i$ are recovered
from $a^2-b^2$ and $ab$.

The relation of these coordinates to the two frames
can be seen without a sign guess.  In the Legendre
family $y^2=x(x-1)(x-t)$, the order-two orbifold
point is $t=1/2$ and the order-three point satisfies
$t^2-t+1=0$.  Transport from the tangential base
point gives the inverse pure conjugating frame
followed by its base-point shear; hence these are
exactly the normalized $W$ and $T$ acting on the
two fiber Tate modules.  The change
$X=4x-2,Y=8y$ at $t=1/2$ gives
$Y^2=X^3-4X$, the twist of $E_4$ by $2$.
At the order-three point, shifting $x$ by
$(1+t)/3$ gives $y^2=x^3+c'^3$ with
$c'=1/\sqrt{-3}$, the twist of $E_3$ by $c'$.
Both polynomial identities are independently
checked in the attached program.

An arithmetic element of multiplier one fixes
$\mathbb Q^{\rm ab}$.  The first twist is then
trivial.  The second twist has character
$\sigma(r_0)/r_0$, since
$(-3)^{1/4}=\zeta_8\sqrt2\,r_0$ and the extra
factors are in $\mathbb Q^{\rm ab}$.
The rational point $(0,1)\in E_3[3]$ spans the
fixed line of its order-three automorphism.
Consequently its frame sign is
$\epsilon_3(T_\sigma)=[\sigma(r_0)\ne r_0]$.
By (E.59), the projective action on $E_4[3]$
gives the identical bit
$\epsilon_2(W_\sigma)=[\sigma(r_0)\ne r_0]$.
Their sum is zero, proving the arithmetic test
against (E.58).

For the other requested tests, the defining
GT$_0$ pair has multiplier one and $f\in F_2'$.
Its automorphism fixes each peripheral generator
up to conjugacy, centralizes the specified
$\mathfrak S_3$ action, and is identity on
$F_2^{\rm ab}$; these follow directly from the
pair description cited above.  Thus none of these
four conditions rejects the designated witness.
The selected maximal AA subgroup is the actual
$G_{\mathbb Q}$ by E.7, so the CM calculation
rejects membership in it.  AA on a subgroup and
membership in that selected subgroup are the
precise tests; no undefined elementwise AA
condition is introduced.  A coherent extension
to five labelled modular points requires the
pentagon, by E.16.  No pentagon verification for
this witness is contained in the supplied
certificate, and the computation (E.58) does not
provide one.

The attachments supply referee reports about the
existence of the designated profinite E.58 element;
they do not supply the original universal relative
peripheral-lifting proof in full.  The computations
above verify the prescribed certificate and all
of the requested tests using that designated
element.  They do not certify that missing
universal lifting argument by a finite test.
Likewise, no conclusion that $\Delta$ vanishes
on GT follows from its vanishing on the
arithmetic subgroup.  The missing value on GT
and the surjectivities in (E.56) are retained as
open assertions, with their first instance
$\widehat\Gamma_{0,5}$ explicit.
\end{proof}

\begin{result}{E.25}{Proved}{A four-puncture automorphism invisible after every filling}
Let $P=\widehat F(a,b,c)$ with fourth peripheral generator
$d=(abc)^{-1}$, and let $q_e:P\twoheadrightarrow
P/\overline{\langle\!\langle e\rangle\!\rangle}$ fill
$e\in\{a,b,c,d\}$.  The product of the four induced
maps on peripheral outer automorphisms has nontrivial
kernel.  In particular, equality of all four filled
tripod outer actions does not determine an outer action
on the four-punctured sphere group.
\end{result}
\begin{proof}
Put $p=ab$ and define a continuous automorphism $T$ by
\[
 T(a)=pap^{-1},\qquad T(b)=pbp^{-1},\qquad T(c)=c.
\]
It fixes $p$ and $d$, and its inverse replaces $p$ by
$p^{-1}$ in the two conjugations.  Thus it preserves
all four peripheral conjugacy classes and induces an
automorphism after every filling.  After filling $a$
or $b$, it is the identity.  After filling $c$, it
is conjugation by $p$ on $\widehat F(a,b)$.  After
filling $d$, the relation $c=p^{-1}$ shows that it is
again conjugation by $p$ on every generator.

Its outer class is nontrivial, with the following finite
certificate.  Evaluate $(a,b,c)$ at
$((12),(23),(13))$ in $\mathfrak S_3$.  With permutation
composition from right to left, the images of $T(b)$
and $T(c)$ are both $(13)$, although the images of
$b,c$ are distinct.  A single conjugation in
$\mathfrak S_3$ cannot have that property.  An inner
representative of $T$ in $P$ would induce just such a
single conjugation after evaluation, a contradiction.
This certificate does not require the kernel of the
evaluation map to be $T$-invariant.

For comparison with (E.53), let
$H=\ker(\widehat F(x,y)\to C_2)$ and write
$\Phi_F(x)=x^\lambda$,
$\Phi_F(y)=f^{-1}y^\lambda f$, with $\lambda=2m+1$.
Let $\delta_2:H\to\widehat F(X,Y)$ have the three
Schreier values $(X,Y,1)$ of E.22.  Equation (E.53)
gives the actual equality
\[
 \delta_2\circ\Phi_F|_H
       =\Phi_F\circ\delta_2.                              \tag{E.60}
\]
On $x^2$ both sides are $X^\lambda$; on $y$ the
extra factor $Y^{c_2(f)}$ commutes with $Y^\lambda$
and cancels.  Finally $\Phi_F(xyx^{-1})$ is a
conjugate in $H$ of $(xyx^{-1})^\lambda$, because
$\lambda$ is odd, so both sides send this generator
to $1$.  Thus the distribution theorem supplies a
genuine covering-and-filling square.  The automorphism
$T$ proves the precise loss of information when one
then retains only the filled outer actions.  It makes
no assertion that the full GT relations cannot imply
(E.51).
\end{proof}

\subsection{Direct genus-one normalizers and the marked surface operad}

\begin{result}{E.26}{Proved}{Two full-completed elliptic stabilizers are self-normalizing}
Use the groups and generators of E.11, and put
$G=\widehat B_4/\langle\omega _4\rangle$,
$Q=\widehat B_3/\langle z^2\rangle$,
$z=(ab)^3$.  Let $H$ be the closure in $G$ of the
nonseparating stabilizer $\langle a,v,d\rangle$ and
let $J$ be the image of the rational-tail subgroup
$\langle a,b\rangle=B_3$.  Both inclusions induce
the full profinite topologies, and
\[
 N_G(H)=H,\qquad N_G(J)=J,\qquad
 C_{\widehat B_3}(a)=\langle a,z\rangle.                 \tag{E.61}
\]
These assertions are independent of the profinite
multitwist and curve-normalizer statements in
\cite[\S5]{ET_Boggi2014}.
\end{result}
\begin{proof}
We first certify all the kernels and topologies used
in the normalizer calculation.  In $B_4$ put
$x=da^{-1}$ and $y=bxb^{-1}$.  Its Artin presentation
gives the split extension
\[
 B_4=F(x,y)\rtimes B_3,
 \quad
 a:(x,y)\longmapsto(x,x^{-1}y),\qquad
 b:(x,y)\longmapsto(y,yx^{-1}y).                          \tag{E.62}
\]
Here is a direct verification.  Write $d=xa$; the
relation $[a,d]=1$ becomes $axa^{-1}=x$.
The relation $bdb=dbd$, together with $aba=bab$,
then gives $aya^{-1}=x^{-1}y$.  Applying $aba=bab$
to $x$ gives $byb^{-1}=yx^{-1}y$.  Conversely these
two automorphisms of $F(x,y)$ satisfy the braid
relation, and $d=xa$ satisfies both remaining
four-strand Artin relations.  The resulting maps
between the two presentations are inverse.
The reduced-word checks, including the calculations
below, are run in \texttt{E\_normalizer\_checks.py}.

Characteristic finite-index subgroups are cofinal
in $F(x,y)$: for each finite index there are finitely
many subgroups, and their intersection is
characteristic.  The corresponding finite quotient
$P$ extends to the finite quotient
$P\rtimes\operatorname{Aut}(P)$ of the semidirect
product.  Thus the kernel in (E.62) has its full
profinite topology and the completed sequence is
split exact.  Under its projection $\omega _4$
maps to $z^2$.  The closed central group generated
by $z^2$ is $\widehat{\mathbb Z}$, since the
exponent-sum map $\widehat B_3\to\widehat{\mathbb Z}$
sends it to $12$ and multiplication by $12$ is
injective.  Quotienting the split extension by
$\langle\omega _4\rangle$ consequently gives
\[
 1\longrightarrow K=\widehat F(X,Y)
 \longrightarrow G\overset q\longrightarrow Q
 \longrightarrow1,
 \qquad X=x,\quad Y=y^{-1}.                              \tag{E.63}
\]
It is the completed Birman sequence, with the
forgetful map $(a,b,d)\mapsto(a,b,a)$.

Set $s=avd$.  The presentation (E.17) becomes
\[
 H_{\rm disc}=\langle a\rangle\times
       \bigl(\langle x\rangle*C_2(s)\bigr),
 \quad q(x)=1,\quad q(s)=-I,\quad q(a)=a.                 \tag{E.64}
\]
Indeed $d=xa$, $v=sx^{-1}a^{-2}$, and the sole
remaining noncommutator relation is $s^2=1$.
Its kernel over $D_{\rm disc}=\langle a,-I\rangle$
has basis $x,sxs$.  Formula (E.62) computes
$sxs=y^{-1}x^{-1}y$.  Hence this kernel is exactly
\[
 L_{\rm disc}=\langle X,YXY^{-1}\rangle
                   \subset F(X,Y).                     \tag{E.65}
\]
In the index-two subgroup defined by the parity of
$Y$, the Schreier basis is
$A=X$, $B=YXY^{-1}$, $C=Y^2$.
Thus $L_{\rm disc}=\langle A,B\rangle$ is a free
factor there.  Retraction killing $C$ proves that
its full profinite topology is induced from $K$.

The quotient $Q/\langle-I\rangle$ is
$\mathrm{PSL}_2(\mathbb Z)^{\wedge}$.
Its normal index-six subgroup obtained by reduction
modulo two is $\widehat F(a^2,b^2)$.  In particular
$a^2$ is a primitive free generator, so
$\langle a\rangle$ has the full procyclic topology.
Moreover $-I$ is independent of this procyclic
group: reduction modulo three puts $a$ in an
order-three subgroup and $-I$ outside it.  Therefore
$\widehat D_{\rm disc}\to Q$ is injective.
Completion is right exact, so the diagram of
$L_{\rm disc}\to H_{\rm disc}\to D_{\rm disc}$
and (E.63), with the two just-proved injections,
proves that $\widehat H_{\rm disc}\to G$ is
injective and that $H\cap K=\widehat L_{\rm disc}$.
This proves the claimed topology and identifies
the actual closed subgroup $H$.

We record an elementary peripheral normalizer
argument.  Let $P$ be the completed fundamental
group of either a sphere with at least three punctures or a
once-punctured torus, and let $I$ be one of its
peripheral procyclic groups.  Then $N_P(I)=I$.
For suppose $g\notin I$ normalizes $I$.  Choose
an open subgroup $U\leq P$ containing $I$ and
excluding $g$.  It corresponds to a connected
finite unramified cover of the punctured surface.
For the sphere it has at least three punctures,
one over each original puncture.  For the torus
its compactification has genus at least one,
since the compactified Riemann--Hurwitz formula
has nonnegative ramification term.  In the latter
case replace $U$ by
\[
 U'=\ker\bigl(U\longrightarrow
      H_1(\overline S_U;\mathbb Z/2)\bigr).
\]
This subgroup contains every cusp inertia of $U$;
the cover of its compactification is unramified
of degree $2^{2g(\overline S_U)}\geq4$.  It
therefore has at least four punctures, still
contains $I$, and still excludes $g$.
Write $V=U$ or $U'$ accordingly.  The double
cosets $VI$ and $VgI$ are distinct, because
$gIg^{-1}=I\subseteq V$ and $g\notin V$.
They describe distinct unramified cusps whose
inertia groups, up to conjugacy in $V$, are
both $I$.  To see this with actual topological
covering maps, choose a discrete representative
$r\in Vg$.  Then $r=hg$ with $h\in V$, so
$rIr^{-1}=hIh^{-1}$; replacing the profinite
coset representative has not changed its
inertia class in $V$.
For a surface with at least three punctures,
the peripheral homology classes have the sole
relation that their sum is zero.  Their distinct
procyclic lines in
$H_1(S_V;\widehat{\mathbb Z})$ are distinct.
The two just-identified inertia groups give the
same line, a contradiction.  This proves the
peripheral assertion.  In particular every
primitive generator of a finite-rank nonabelian free profinite
group has its own procyclic group as normalizer
and centralizer.

Now $N_K(L)=L$ for
$L=\overline{\langle X,YXY^{-1}\rangle}$.
First work in $\widehat F(A,B,C)$ as above.
If $h$ normalizes its free factor $\widehat F(A,B)$,
let $r$ be the retraction killing $C$.
The element $r(h)^{-1}h$ centralizes the free
factor: on this factor both conjugations have
the same image under $r$.  It centralizes $A,B$,
so the primitive-generator centralizer assertion
just proved forces it to be $1$.
Thus $N_{\widehat F(A,B,C)}(L)=L$.
An odd normalizer in $K$ would have the form
$hY$ and would imply
\[
 h\langle B,CAC^{-1}\rangle h^{-1}=\langle A,B\rangle.
\]
Kill $B$ in this equality.  In $\widehat F(A,C)$,
the image of $h$ times $C$ normalizes
$\langle A\rangle$, so the $C$-coordinate of
$h$ in abelianization is $-1$.
Instead kill $A$.  In $\widehat F(B,C)$ the
image of $h$ normalizes $\langle B\rangle$, so
the same coordinate is zero.  This contradiction
excludes odd normalizers.

Also $N_Q(D)=D$ for $D=\langle a,-I\rangle$.
In the quotient $\mathrm{PSL}_2(\mathbb Z)^{\wedge}$,
the image of a normalizer of $\langle a\rangle$
in $\mathfrak S_3$ normalizes the transposition
subgroup, hence lies in it.  Multiply by a
suitable power of $a$ to enter the index-six
free subgroup.  The remaining element normalizes
$\langle a^2\rangle$, whose primitive-generator
normalizer is itself.  Lifting back to $Q$
proves the assertion.  Finally a normalizer of
$H$ projects into $D$, so multiplying it by an
element of $H$ reduces it to $K$.  It then
normalizes $H\cap K=L$, and hence belongs to
$L$.  This proves $N_G(H)=H$.

For $J$, the discrete cutting sequence is
$1\to\langle z^2\rangle\to B_3\to
\mathrm{SL}_2(\mathbb Z)\to1$.
In (E.62) one has
$\omega _4=([X,Y],z^2)$, so in (E.63)
$z^2=[X,Y]^{-1}$.  This identifies its kernel
with the actual cusp group
$I=\langle[X,Y]\rangle\subset K$.
That cyclic group has the full topology: in the
index-two free subgroup used above,
$[X,Y]=AB^{-1}$ is primitive.  The right-exact
completion diagram proves that
$\widehat B_3\to G$ is injective and that
$J\cap K=I$.  A normalizer of $J$ may be
multiplied by an element of $J$ to enter $K$,
since $q(J)=Q$.  It then normalizes $I$, so
the finite-cover argument puts it in $I$.
This proves $N_G(J)=J$.
Notice that $Z(J)=\langle z\rangle$, whereas
the node-twist group is $\langle z^2\rangle$;
their index-two difference has not been suppressed.

The same quotient by $\langle z\rangle$ maps
$\widehat B_3$ onto
$\mathrm{PSL}_2(\mathbb Z)^{\wedge}$, with
injected procyclic kernel detected by exponent
sum.  The preceding primitive-generator argument
shows that the centralizer of the image of $a$
is $\langle a\rangle$.  Its inverse image is
$\langle a,z\rangle$, an abelian group.  Hence
it is exactly $C_{\widehat B_3}(a)$.
All three equalities in (E.61) are proved.
\end{proof}

\begin{result}{E.27}{Proved}{Genus zero detects strict marked-pants operad automorphisms}
Let $\mathbf S$ be the modular operad of surfaces
with marked pants objects and parametrized labelled
boundaries in \cite[Definitions 3.1--3.3]{ET_BonattoRobertson2026},
and complete its morphism groups and torsors fully.
For continuous strict automorphisms fixing every
object, restriction is injective:
\[
 \operatorname{Aut}_0(\widehat{\mathbf S})
   \lhook\joinrel\longrightarrow
 \operatorname{Aut}_0(\operatorname{tr}_{\leq0}
                     \widehat{\mathbf S})
   =\widehat{\operatorname{GT}}.                          \tag{E.66}
\]
The map is restriction along the actual genus-zero
cyclic suboperad; the modular operations here are
gluing, contraction, and relabelling.  No
forgetful operation is assumed in this statement.
\end{result}
\begin{proof}
Let $F$ be in the kernel.  The genus-zero
identification in (E.66) is
\cite[Propositions 3.16 and 4.7]{ET_BonattoRobertson2026}
and the parenthesized ribbon-braid description.
We prove directly that $F$ fixes the remaining
generator, the one-holed-torus move
$\sigma=aba\in B_3$.

Cap the two external boundary circles in a
chosen genus-one, two-boundary object by marked
disks.  Algebraically this quotients its mapping
class group by the two central boundary twists.
Since these twists are compositions with the
genus-zero annulus generator, $F$ fixes them
and induces an automorphism $\alpha$ of
$G=\widehat\Gamma_{1,2}=\widehat B_4/
\langle\omega _4\rangle$.
By \cite[Theorem 4.3(i) and Corollary C]{ET_MN2022},
there are $g\in G$ and a unique
$A=(\lambda,f)\in\widehat{\operatorname{GT}}$ such that
\[
             \alpha=\operatorname{Inn}(g)\phi_A,         \tag{E.67}
\]
where $\phi_A$ is the genuine four-strand action.

Use the nonseparating cut defining $H$ in E.26.
The ordered cut surface is a four-boundary
sphere.  Its mapping classes, followed by
contraction of the two specified boundaries,
give the branch-preserving subgroup $H^+$.
The two remaining boundary twists disappear
on capping, and the two sewn-boundary twists
give $\langle a\rangle$.  Consequently its image
in
\[
 R=H/\langle a\rangle
    =\langle v,d\mid(vd)^2=1\rangle^{\wedge}
\]
is exactly
$K_{\rm geo}=\ker(v\mapsto1,d\mapsto0)=
\widehat F(d,v^2)$.
These image equalities follow first from cutting
and then from compactness of the completed
source; E.26 verifies the topology of the target.
To include the branch-swapping coset, take a
cut-surface morphism whose target has its two
sewing labels exchanged and apply contraction.
The contracted source and target graph coincide.
Thus compatibility with relabelling and
contraction makes $F$ preserve the full $H$,
while its identity genus-zero restriction fixes
the image $K_{\rm geo}$ pointwise.

E.11 proves that $\phi_A$ also preserves $H$.
Thus (E.67) and $N_G(H)=H$ imply $g\in H$.
Modulo $\langle a\rangle$, write $\bar g$ for
its image.  The action on $R$ is
\[
 \phi_A(v)=v^\lambda,
 \qquad \phi_A(d)=h^{-1}d^\lambda h,
 \qquad h=f(v,d^2).                                     \tag{E.68}
\]
The element $h$ is in both index-two kernels,
since $f$ has zero abelianization.  On the
abelianization of $K_{\rm geo}$, (E.68) therefore
acts by $\lambda$ times the identity.
Conjugation by an element of its odd coset
instead has, in the basis $(d,v^2)$, matrix
\[
                 \begin{pmatrix}-1&0\\-1&1\end{pmatrix}.
                                                               \tag{E.69}
\]
Indeed $vdv^{-1}=d^{-1}v^{-2}$ and $v$ fixes
$v^2$.  The off-diagonal entry is nonzero
already modulo two.  The equality
$\operatorname{Inn}(\bar g)\phi_A|_{K_{\rm geo}}=1$
therefore forces $\bar g\in K_{\rm geo}$
and $\lambda=1$.

The centralizer of $K_{\rm geo}$ in $R$ is
trivial.  Its even part is the center of the
nonabelian free profinite group, which is
trivial by the primitive-generator centralizer
argument in E.26.  An odd element cannot
centralize it, by (E.69).
An automorphism of $R$ fixing this normal
subgroup pointwise must consequently fix every
element $r$: the difference between its value
on $r$ and $r$ centralizes $K_{\rm geo}$.
It follows that
$\phi_A=\operatorname{Inn}(\bar g^{-1})$ on all $R$.

Now use the other index-two subgroup
$K_{\rm std}=\ker(v\mapsto0,d\mapsto1)=
\widehat F(v,d^2)$.
Conjugation by its odd coset has the identical
nonscalar matrix (E.69), now in the basis
$(v,d^2)$; the calculation is
$dvd^{-1}=v^{-1}d^{-2}$.
Since (E.68) acts identically on this
abelianization, $\bar g$ must also belong to
$K_{\rm std}$.
Thus the standard free-group action of
$(1,f)$ on the basis $(v,d^2)$ is inner.
For clarity its faithfulness can be checked
without an outer-automorphism classification.
An inner representative fixing $v$ is conjugation
by $v^t$, because $C_{K_{\rm std}}(v)=\langle v\rangle$.
Its equality with the second formula in (E.68)
implies $hv^t\in\langle d^2\rangle$.
Abelianization, and the fact that $h$ is a
commutator word, give $t=0$ and this remaining
power equal to zero.  Hence $h=1$, and
evaluation on the free basis $(v,d^2)$ gives
$f=1$.  Therefore $A=1$ and $\alpha$ is inner.

It remains to remove that inner ambiguity
without using a forgetting map in $\mathbf S$.
Glue the one-holed torus to a pair of pants
along its boundary and cap the other two
boundary circles by marked disks.  The
resulting genuine homomorphism is exactly
\[
        j:\widehat B_3\longrightarrow G,
                 (a,b)\longmapsto(a,b),
\]
with image $J$ of E.26.  It is injective by
that result.  The strict operad automorphism
commutes with this gluing; consequently
$\operatorname{Inn}(g)$ preserves $J$ and
$g\in J$ by $N_G(J)=J$.
Its action on $\widehat B_3$ is therefore inner.
The twist $a$ on the one-holed torus is itself
the contraction of a genus-zero boundary twist,
so $F(a)=a$.
By $C_{\widehat B_3}(a)=\langle a,z\rangle$
there is $\nu\in\widehat{\mathbb Z}$ such that
\[
                   F(\sigma)=a^{-\nu}\sigma a^\nu.      \tag{E.70}
\]
The central factor $z$ has no effect on this
conjugation.

Finally apply the actual genus-one relation
(G2) of \cite[Theorem 3.18]{ET_BonattoRobertson2026}.
Its calculation for (E.70) is the one displayed
in that source, (6.3)--(6.5), pp.~54--55.
Here it can be checked in the four-strand
presentation itself.  Put
$U=dbd$ and $\delta=Ud^{-1}aU$.
This is the Garside half twist, and
$\delta d\delta^{-1}=a$.
Since all the genus-zero factors of (G2) are
fixed, substitution of (E.70) in its two
one-holed-torus occurrences gives
\[
              \delta=d^{-\nu}\delta d^\nu,
       \qquad 1=d^{-\nu}a^\nu=X^{-\nu}.                 \tag{E.71}
\]
The equality of the two expressions for
$\delta$ and its conjugation formula are
verified in the faithful rank-four Artin action
by \texttt{E\_normalizer\_checks.py}.
The last equality of (E.71) uses $X=da^{-1}$
from (E.63), and $[a,d]=1$.
Since $X$ is a free generator of $K$, its
abelianization forces $\nu=0$.
Thus $F(\sigma)=\sigma$.

Every discrete morphism of $\mathbf S$ is
generated by $\tau,\beta,\alpha,\sigma$ using
composition, gluing, contractions, and
relabelling, by
\cite[Proposition 3.8 and its proof]{ET_BonattoRobertson2026}.
The first three lie in genus zero, and the
fourth has just been fixed.  These morphisms
are dense in every completed morphism torsor.
Continuity now gives $F=1$, proving (E.66).

E.29 supplies the full completed cutting groups
from the exit tower by using the diagonal
stabilizer of each pinned crossing, with both
endpoints retained.  Recovering all crossing-coset
maps from the resulting actions remains the
faithfulness obligation recorded in E.23.
The theorem above
settles the kernel for the explicitly specified
operad itself.
\end{proof}

\begin{result}{E.28}{Proved}{The two elliptic readings give an involution of GT}
In $R=\langle v,d\mid(vd)^2=1\rangle^{\wedge}$,
write $D=d$, $V=v^2$, and let $c_2(f)$ have the
Schreier meaning of E.22.  For
$F=(\lambda,f)\in\widehat{\operatorname{GT}}$, put
\[
 \mathcal J(F)=\left(\lambda,
       V^{2c_2(f)}f(v,d^2)^{-1}D^{4c_2(f)}\right),        \tag{E.72}
\]
where the second entry is read in the free basis
$(D,V)$ of $K_{\rm geo}$.
This pair belongs to $\widehat{\operatorname{GT}}$,
and $\mathcal J$ is a continuous involutive
automorphism of that group.  Its fixed-point
equation is exactly (E.50), equivalently (E.51).
It fixes the arithmetic subgroup $G_{\mathbb Q}$.
The assertion that $\mathcal J$ is the identity
on all of GT remains the explicit elliptic
comparison of E.23.
\end{result}
\begin{proof}
The point needing proof is that the pair in
(E.72) satisfies the five-point relations.  We
construct its action on the actual five-point
group, using E.9, rather than infer those
relations from a three-point action.

Retain $D_4=C_2(x_1)\amalg\cdots\amalg C_2(x_4)$
and $K_4$ of E.9.  Put
\[
 p=x_1x_2,\qquad
 B=\overline{\langle p,x_3,x_4\rangle},\qquad
 L=B\cap K_4.
\]
The index-two subgroup of $D_4$ defined by
$x_1,x_2\mapsto1$, $x_3,x_4\mapsto0$ has the
explicit free-product decomposition
\[
 \langle p\rangle\amalg C_2(x_3)\amalg C_2(x_4)
 \amalg C_2(x_1x_3x_1)\amalg C_2(x_1x_4x_1).
                                                               \tag{E.73}
\]
Reidemeister--Schreier with representatives
$1,x_1$ gives exactly these factors and only
their displayed order-two relations.
Thus $B$ is a free factor of this open subgroup
and has the full induced topology.  Its
index-two subgroup $L$ is free on
\[
             p,\qquad q_0=x_3px_3,
                  \qquad r=x_3x_4.                     \tag{E.74}
\]
It also has the full topology in $K_4$.

The disk double cover defining E.9 gives a
geometric identification of this subgroup.
Remove a small disk around branch points
$x_1,x_2$ before taking the cover.  Its
boundary has monodromy $p$, which is even;
the remaining branch points are $x_3,x_4$.
The double cover is a genus-zero surface
with four boundary components.  It is the
complement of an annular neighborhood of the
nonseparating lift of the arc joining
$x_1,x_2$.  Hence $L$ is precisely the
point-pushing kernel for that cut surface.
The two old punctures have representatives
\[
                 c=pr,
       \qquad c'=x_3cx_3=q_0r^{-1}.                     \tag{E.75}
\]
The other two boundary classes are $p,q_0$,
with their boundary orientations reversed
when viewed as punctures of the cut surface.

Let $\widetilde H=\langle a,v,d\rangle
\subset\widehat B_4$ be the preimage of $H$
before the central quotient; here
$a=\sigma_1$, $v=\sigma_2\sigma_1^2\sigma_2$,
$d=\sigma_3$.  In the extension of E.9 the
cut stabilizer is exactly
\[
 H_3=(L\rtimes\widetilde H)/
               \langle(c,\omega _4)\rangle.        \tag{E.76}
\]
The discrete equality follows by cutting and
forgetting the third point: its kernel is
the fundamental group of the cut surface
just described, and every mapping class
preserving the cut can be lifted to fix the
third point in its complement.  Its quotient
is the discrete $H$ of E.26.
The full-topology equality follows from the
completed discrete extension diagram: $L$
injects into $K_4$ by (E.73), and
$\widehat H$ injects into $G$ by E.26;
right exactness and the two injections give
injectivity for the middle group.
The central relation in (E.76) lies in the
displayed subgroup because $c=px_3x_4\in L$.

Its branch-preserving index-two subgroup
$H_3^+$ is the inverse image of
$\ker(v\mapsto1,d\mapsto0)$ in $H$.
Cutting with both branches ordered gives
the exact quotient
\[
 H_3^+/\langle a\rangle=\widehat\Gamma_{0,5},
 \qquad
 \widehat\Gamma_{0,5}\longrightarrow
        H^+/\langle a\rangle
          =K_{\rm geo}=\widehat\Gamma_{0,4}.             \tag{E.77}
\]
The second map forgets the third, moving
point; its kernel is $L$.
These are quotients of the actual completed
cutting sequences.  Thus the five-point
group in (E.77) carries its full topology.

The genuine GT action of E.9 preserves
(E.76).  To check the only additional
subgroup, apply the first-strand doubling
homomorphism
$A_4\to\widehat B_5$ of
\cite[\S0.3, pp.~7--8]{ET_HS2000}.
In its subgroup fixing the last strand,
the three point-pushing meridians map exactly as
\[
 x_{14}\longmapsto x_{15}x_{25}=p,
 \quad x_{24}\longmapsto x_{35}=x_3,
 \quad x_{34}\longmapsto x_{45}=x_4.
\]
For example the first image is
$\sigma_4\sigma_3\sigma_2\sigma_1^2
\sigma_2\sigma_3^{-1}\sigma_4^{-1}$,
which is the product $x_{15}x_{25}$ by
cancelling the middle inverse strings.
The genuine GT action commutes
with this doubling, and sends the two latter
meridians to conjugates of their odd unit
powers.  It therefore descends after imposing
$x_3^2=x_4^2=1$ and preserves $B$.
It preserves parity and hence $L$.
It preserves $\widetilde H$ by E.11 and sends
both $c$ and $\omega_4$ to their
$\lambda$th powers.  It consequently preserves
the exact central relation in (E.76), the
branch character, and $\langle a\rangle$.
This constructs an actual automorphism
$\beta_F$ of $\widehat\Gamma_{0,5}$ compatible
with the quotient (E.77).

By \cite[Corollary 2.8(i)]{ET_HMM2022}, write
its outer class as an element of
$\widehat{\operatorname{GT}}\times\mathfrak S_5$.
We check that its permutation factor is
trivial.  It preserves the kernel of
forgetting the moving point, so its
$\mathfrak S_5$ factor fixes that label.
Indeed the five forgetting kernels are
distinct: a pair twist involving one
specified label and a third label is killed
on forgetting the former but survives on
forgetting a different label.
On the quotient $K_{\rm geo}$, (E.68)
preserves all three peripheral classes
$d,v^2,vdv^{-1}$ with exponent $\lambda$.
For the last one use that $\lambda$ is odd
and $h=f(v,d^2)$ lies in the even kernel.
Thus the remaining $\mathfrak S_4$ factor
acts trivially on the three partitions into
two pairs, and lies in its Klein four group.
Finally (E.75) distinguishes the two original
labels individually: the genuine action
sends $c$ to $c^\lambda$, and its value on
$x_3$ is an odd element of $B$.
Consequently its value on $c'$ is conjugate
to $(c')^\lambda$ by the even element
$F(x_3)x_3\in L$.
The permutation factor fixes both of these
labels.  Every nonidentity element of the
Klein four group moves all four letters,
so this factor is the identity.
The outer class of $\beta_F$ therefore
belongs to the GT factor itself.

We compute its parameter from (E.77).
The Schreier exponent calculation of E.22,
with the roles of $v,d$ exchanged, gives
\[
       [h]_{K_{\rm geo}}=(4c_2(f),2c_2(f))
                     \quad\text{in the basis }(D,V).
\]
Conjugating (E.68) by
$D^{-4c_2(f)}h$ fixes $D^\lambda$ and sends
$V$ to $f_{\rm geo}^{-1}V^\lambda f_{\rm geo}$,
where
$f_{\rm geo}=V^{2c_2(f)}h^{-1}D^{4c_2(f)}$.
This word has zero abelianization.  It is
therefore exactly the normalized pair
parameter of the GT factor just obtained.
This proves membership asserted in (E.72),
including the five-point relation.
The construction by genuine actions,
restriction, quotient, and projection to
the GT factor also proves continuity and
the homomorphism law for $\mathcal J$.

For involutivity let $S$ be the automorphism
of $R$ exchanging $v,d$, and let
$\iota:\widehat{\operatorname{GT}}\to
\operatorname{Out}(R)$ be the standard action
(E.68).
This map is injective: an inner representative
must lie in $K_{\rm std}$ by the nonscalar
odd-coset matrix (E.69), and the free-basis
faithfulness argument in E.27 then applies.
By the construction of $f_{\rm geo}$,
$\phi_{\mathcal J(F)}$ and
$S\phi_FS$ agree, up to an inner conjugation,
on $K_{\rm std}$.
Their equality on that normal subgroup
extends to $R$, since its centralizer is
trivial, as proved in E.27.
Hence
\[
       \iota(\mathcal J(F))=S\iota(F)S,
                  \qquad\mathcal J^2(F)=F.              \tag{E.78}
\]
This proves that $\mathcal J$ is an
involutive automorphism.
Equality of its second component with
$f(D,V)$ is precisely (E.50), by multiplication
by $D^{-4c_2(f)}h$; E.22 proves its equivalent
full-braid form (E.51).
For arithmetic $F$, the Galois action
commutes with the two specified geometric
boundary maps by E.4, and the orbifold
parameter is the same $F$ by E.12.
Thus (E.50) holds for these elements,
giving the last assertion.
\end{proof}
